% !TeX program = lualatex
% ALIUS Bulletin native visual reconstruction source.
% Generated from off-repo reference layout data; does not include or input a pre-existing article PDF.
\ifdefined\ALIUSIssueBuild
\else
\documentclass{article}
\usepackage[paperwidth=595.0000bp,paperheight=842.0000bp,margin=0bp]{geometry}
\usepackage{graphicx}
\usepackage{xcolor}
\usepackage{tikz}
\usepackage{iftex}
\usepackage[hidelinks]{hyperref}
\ifPDFTeX
  \usepackage[T1]{fontenc}
  \usepackage[utf8]{inputenc}
  \usepackage{textcomp}
  \usepackage{amssymb}
\else
  \usepackage{fontspec}
\fi
\pagestyle{empty}
\fi
\ifPDFTeX
  \providecommand{\ALIUSDeclareUnicodeFallbacks}{%
    \DeclareUnicodeCharacter{00AC}{\ensuremath{\neg}}%
    \DeclareUnicodeCharacter{00BE}{\ensuremath{\frac{3}{4}}}%
    \DeclareUnicodeCharacter{0101}{\=a}%
    \DeclareUnicodeCharacter{0107}{\'c}%
    \DeclareUnicodeCharacter{012B}{\=\i}%
    \DeclareUnicodeCharacter{0131}{\i}%
    \DeclareUnicodeCharacter{0142}{\l}%
    \DeclareUnicodeCharacter{015B}{\'s}%
    \DeclareUnicodeCharacter{016B}{\=u}%
    \DeclareUnicodeCharacter{025C}{q}%
    \DeclareUnicodeCharacter{0278}{t}%
    \DeclareUnicodeCharacter{02AE}{Th}%
    \DeclareUnicodeCharacter{02B0}{ff}%
    \DeclareUnicodeCharacter{02B4}{ffi}%
    \DeclareUnicodeCharacter{02BE}{ft}%
    \DeclareUnicodeCharacter{02BF}{fi}%
    \DeclareUnicodeCharacter{02C0}{fl}%
    \DeclareUnicodeCharacter{02C1}{fi}%
    \DeclareUnicodeCharacter{02D9}{Th}%
    \DeclareUnicodeCharacter{02EF}{gy}%
    \DeclareUnicodeCharacter{0394}{\ensuremath{\Delta}}%
    \DeclareUnicodeCharacter{03B2}{\ensuremath{\beta}}%
    \DeclareUnicodeCharacter{03BA}{\ensuremath{\kappa}}%
    \DeclareUnicodeCharacter{068D}{-}%
    \DeclareUnicodeCharacter{08BC}{ti}%
    \DeclareUnicodeCharacter{099E}{ti}%
    \DeclareUnicodeCharacter{1E43}{\d{m}}%
    \DeclareUnicodeCharacter{1E47}{\d{n}}%
    \DeclareUnicodeCharacter{1E63}{\d{s}}%
    \DeclareUnicodeCharacter{201C}{``}%
    \DeclareUnicodeCharacter{201D}{''}%
    \DeclareUnicodeCharacter{2260}{\ensuremath{\neq}}%
    \DeclareUnicodeCharacter{25A1}{\ensuremath{\square}}%
  }%
  \ifdefined\ALIUSIssueBuild\else\ALIUSDeclareUnicodeFallbacks\fi
  \providecommand{\ALIUSFontLatoLight}{\fontfamily{lato-LF}\fontseries{l}\selectfont}
  \providecommand{\ALIUSFontLatoRegular}{\fontfamily{lato-LF}\fontseries{m}\selectfont}
  \providecommand{\ALIUSFontLatoItalic}{\fontfamily{lato-LF}\fontseries{m}\itshape}
  \providecommand{\ALIUSFontLatoLightItalic}{\fontfamily{lato-LF}\fontseries{l}\itshape}
  \providecommand{\ALIUSFontCormorantRegular}{\fontfamily{CormorantGaramond-LF}\fontseries{m}\selectfont}
  \providecommand{\ALIUSFontCormorantLight}{\fontfamily{CormorantGaramond-LF}\fontseries{l}\selectfont}
  \providecommand{\ALIUSFontCormorantMedium}{\fontfamily{CormorantGaramond-LF}\fontseries{medium}\selectfont}
  \providecommand{\ALIUSFontCormorantBold}{\fontfamily{CormorantGaramond-LF}\fontseries{b}\selectfont}
  \providecommand{\ALIUSFontCormorantItalic}{\fontfamily{CormorantGaramond-LF}\fontseries{m}\itshape}
  \providecommand{\ALIUSFontTimes}{\fontfamily{ptm}\selectfont}
  \providecommand{\ALIUSFontCalibri}{\fontfamily{phv}\selectfont}
  \providecommand{\ALIUSFontCambria}{\fontfamily{ptm}\selectfont}
  \providecommand{\ALIUSFontSymbolFallback}{\fontfamily{ptm}\selectfont}
\else
  \providecommand{\ALIUSFontLatoLight}{\fontspec{Lato Light}}
  \providecommand{\ALIUSFontLatoRegular}{\fontspec{Lato Regular}}
  \providecommand{\ALIUSFontLatoItalic}{\fontspec{Lato Italic}}
  \providecommand{\ALIUSFontLatoLightItalic}{\fontspec{Lato Light Italic}}
  \providecommand{\ALIUSFontCormorantRegular}{\fontspec{Cormorant Garamond Regular}}
  \providecommand{\ALIUSFontCormorantLight}{\fontspec{Cormorant Garamond Light}}
  \providecommand{\ALIUSFontCormorantMedium}{\fontspec{Cormorant Garamond Medium}}
  \providecommand{\ALIUSFontCormorantBold}{\fontspec{Cormorant Garamond Bold}}
  \providecommand{\ALIUSFontCormorantItalic}{\fontspec{Cormorant Garamond Italic}}
  \providecommand{\ALIUSFontTimes}{\fontspec{Times New Roman}}
  \providecommand{\ALIUSFontCalibri}{\fontspec{Calibri}}
  \providecommand{\ALIUSFontCambria}{\fontspec{Cambria}}
  \providecommand{\ALIUSFontSymbolFallback}{\fontspec{Times New Roman}}
\fi
\providecommand{\ALIUSPullQuoteOpen}{\textquotedblleft}
\providecommand{\ALIUSPullQuoteClose}{\textquotedblright}
\definecolor{ALIUSC000000}{HTML}{000000}
\definecolor{ALIUSC0000FF}{HTML}{0000FF}
\definecolor{ALIUSC1F8135}{HTML}{1F8135}
\definecolor{ALIUSC595959}{HTML}{595959}
\definecolor{ALIUSC767171}{HTML}{767171}
\definecolor{ALIUSCBFBFBF}{HTML}{BFBFBF}
\providecommand{\ALIUSPlacedText}[6]{%
  \node[anchor=north west,inner sep=0pt,outer sep=0pt,text=#4] at (#1bp,-#2bp) {%
    \resizebox{#3bp}{!}{{#5\fontsize{#6bp}{#6bp}\selectfont #6\kern0pt}}%
  };%
}
\providecommand{\ALIUSPlacedTextContent}[7]{%
  \node[anchor=north west,inner sep=0pt,outer sep=0pt,text=#4] at (#1bp,-#2bp) {%
    \resizebox{#3bp}{!}{{#5\fontsize{#6bp}{#6bp}\selectfont #7}}%
  };%
}
\ifdefined\ALIUSIssueBuild
\else
\begin{document}
\fi
% Page 1
\thispagestyle{empty}
\noindent\begin{tikzpicture}[remember picture,overlay,x=1bp,y=1bp,shift={(current page.north west)}]
  \fill[ALIUSC1F8135] (371.520bp,-172.320bp) rectangle ++(115.200bp,-0.240bp);
  \fill[ALIUSC1F8135] (371.520bp,-253.920bp) rectangle ++(74.640bp,-0.240bp);
  \fill[ALIUSC1F8135] (371.520bp,-310.320bp) rectangle ++(107.760bp,-0.240bp);
  \fill[ALIUSCBFBFBF] (365.760bp,-149.280bp) rectangle ++(0.240bp,-92.400bp);
  \fill[ALIUSC0000FF] (87.840bp,-314.640bp) rectangle ++(159.840bp,-0.480bp);
  \fill[ALIUSCBFBFBF] (365.760bp,-241.680bp) rectangle ++(0.240bp,-2.160bp);
  \fill[ALIUSCBFBFBF] (365.760bp,-243.840bp) rectangle ++(0.240bp,-99.840bp);
  \ALIUSPlacedTextContent{90.000}{72.019}{306.736}{ALIUSC000000}{\ALIUSFontLatoRegular}{19.920}{Indian philosophy and the value of}
  \ALIUSPlacedTextContent{90.000}{96.019}{244.156}{ALIUSC000000}{\ALIUSFontLatoRegular}{19.920}{transformative experiences}
  \ALIUSPlacedTextContent{371.520}{149.184}{82.102}{ALIUSC000000}{\ALIUSFontLatoRegular}{11.040}{Monima Chadha}
  \ALIUSPlacedTextContent{371.520}{162.359}{115.270}{ALIUSC1F8135}{\ALIUSFontLatoLight}{9.120}{monima.chadha@monash.edu}
  \ALIUSPlacedTextContent{371.520}{173.159}{124.131}{ALIUSC000000}{\ALIUSFontLatoLight}{9.120}{Department of Philosophy and}
  \ALIUSPlacedTextContent{371.520}{183.959}{138.693}{ALIUSC000000}{\ALIUSFontLatoLight}{9.120}{Monash Centre for Consciousness}
  \ALIUSPlacedTextContent{214.230}{187.687}{134.770}{ALIUSC000000}{\ALIUSFontLatoLight}{18.960}{Monima Chadha}
  \ALIUSPlacedTextContent{87.840}{190.529}{126.427}{ALIUSC000000}{\ALIUSFontLatoLight}{16.080}{An interview with}
  \ALIUSPlacedTextContent{371.520}{194.759}{107.611}{ALIUSC000000}{\ALIUSFontLatoLight}{9.120}{and Contemplative studies,}
  \ALIUSPlacedTextContent{371.520}{205.559}{113.556}{ALIUSC000000}{\ALIUSFontLatoLight}{9.120}{Monash University, Australia}
  \ALIUSPlacedTextContent{87.840}{226.008}{222.921}{ALIUSC000000}{\ALIUSFontLatoLight}{12.960}{by Edvard Avilés and Matthieu Koroma}
  \ALIUSPlacedTextContent{371.520}{230.784}{68.390}{ALIUSC000000}{\ALIUSFontLatoRegular}{11.040}{Edvard Aviles}
  \ALIUSPlacedTextContent{371.520}{243.959}{74.753}{ALIUSC1F8135}{\ALIUSFontLatoLight}{9.120}{eda39@cornell.edu}
  \ALIUSPlacedTextContent{371.520}{254.759}{109.250}{ALIUSC000000}{\ALIUSFontLatoLight}{9.120}{Sage School of Philosophy,}
  % ALIUS normalized citation panel begin
  % ALIUS normalized citation panel text: Cite as: Chadha, M., Avilés, E. \& Koroma, M. (2021). Indian philosophy and the value of transformative experiences. An interview with Monima Chadha by Edvard Aviles and Matthieu Koroma. ALIUS Bulletin, 5, 29-38,
  \fill[white] (86.340bp,-253.731bp) rectangle ++(276.680bp,-95.500bp);
  \node[anchor=north west,inner sep=0pt,outer sep=0pt,text=ALIUSC000000] at (87.840bp,-255.731bp) {\begin{minipage}{273.680bp}{\ALIUSFontLatoLight\fontsize{9.200bp}{10.700bp}\selectfont\raggedright Cite as: Chadha, M., Avilés, E. \& Koroma, M. (2021). Indian philosophy and the value of transformative experiences. An interview with Monima Chadha by Edvard Aviles and Matthieu Koroma. ALIUS Bulletin, 5, 29-38,\par}\end{minipage}};
  \node[anchor=north west,inner sep=0pt,outer sep=0pt,text=ALIUSC1F8135] at (87.840bp,-303.031bp) {{\ALIUSFontLatoLight\fontsize{8.600bp}{8.600bp}\selectfont\href{https://doi.org/10.34700/rn10-az29}{https://doi.org/10.34700/rn10-az29}}};
  % ALIUS normalized citation panel end
  \ALIUSPlacedTextContent{371.520}{265.559}{96.254}{ALIUSC000000}{\ALIUSFontLatoLight}{9.120}{Cornell University, USA}
  \ALIUSPlacedTextContent{371.520}{287.184}{88.206}{ALIUSC000000}{\ALIUSFontLatoRegular}{11.040}{Matthieu Koroma}
  \ALIUSPlacedTextContent{371.520}{300.359}{107.740}{ALIUSC1F8135}{\ALIUSFontLatoLight}{9.120}{matthieu.koroma@uliege.be}
  \ALIUSPlacedTextContent{371.520}{311.159}{110.514}{ALIUSC000000}{\ALIUSFontLatoLight}{9.120}{Physiology of Cognition lab,}
  \ALIUSPlacedTextContent{371.520}{321.959}{88.024}{ALIUSC000000}{\ALIUSFontLatoLight}{9.120}{GIGA Consciousness,}
  \ALIUSPlacedTextContent{371.520}{332.759}{108.490}{ALIUSC000000}{\ALIUSFontLatoLight}{9.120}{University of Liege, Belgium}
  \ALIUSPlacedTextContent{90.000}{375.781}{55.719}{ALIUSC000000}{\ALIUSFontLatoRegular}{13.920}{Abstract}
  \ALIUSPlacedTextContent{90.000}{404.316}{418.265}{ALIUSC000000}{\ALIUSFontLatoLight}{12.000}{In this interview, we engage in a cross-cultural discussion about the diversity of}
  \ALIUSPlacedTextContent{90.000}{418.716}{418.264}{ALIUSC000000}{\ALIUSFontLatoLight}{12.000}{consciousness. Indian philosophy can seem quite cryptic and difficult to follow}
  \ALIUSPlacedTextContent{90.000}{433.116}{418.264}{ALIUSC000000}{\ALIUSFontLatoLight}{12.000}{because it is a primary oral tradition. However, Monima Chadha has developed a}
  \ALIUSPlacedTextContent{90.000}{447.516}{418.264}{ALIUSC000000}{\ALIUSFontLatoLight}{12.000}{series of work aiming at introducing the rich insights of Indian philosophy of mind}
  \ALIUSPlacedTextContent{90.000}{461.916}{418.264}{ALIUSC000000}{\ALIUSFontLatoLight}{12.000}{into the western literature. Key aspects of metaphysical debates such as the}
  \ALIUSPlacedTextContent{90.000}{476.316}{418.267}{ALIUSC000000}{\ALIUSFontLatoLight}{12.000}{nature of consciousness and the notion of no-self are framed in Indian philosophy}
  \ALIUSPlacedTextContent{90.000}{490.716}{418.265}{ALIUSC000000}{\ALIUSFontLatoLight}{12.000}{as primarily guided by the aim to achieve enlightenment or liberation. Altered}
  \ALIUSPlacedTextContent{90.000}{505.116}{418.265}{ALIUSC000000}{\ALIUSFontLatoLight}{12.000}{states of consciousness are considered as transformative experiences and}
  \ALIUSPlacedTextContent{90.000}{519.516}{418.265}{ALIUSC000000}{\ALIUSFontLatoLight}{12.000}{meditation is highlighted as an additional resource in the Indian tradition to gain}
  \ALIUSPlacedTextContent{90.000}{533.916}{214.752}{ALIUSC000000}{\ALIUSFontLatoLight}{12.000}{insight into the nature of consciousness.}
  \ALIUSPlacedTextContent{90.000}{560.316}{47.340}{ALIUSC000000}{\ALIUSFontLatoItalic}{12.000}{keywords}
  \ALIUSPlacedTextContent{137.339}{560.316}{366.890}{ALIUSC000000}{\ALIUSFontLatoLightItalic}{12.000}{: Indian philosophy, consciousness, meditation, transformative experiences,}
  \ALIUSPlacedTextContent{90.000}{574.716}{102.549}{ALIUSC000000}{\ALIUSFontLatoLightItalic}{12.000}{psychedelics, no-self}
  \ALIUSPlacedTextContent{90.000}{618.463}{418.395}{ALIUSC1F8135}{\ALIUSFontLatoLight}{12.357}{You are an Associate Professor in Philosophy at Monash University}
  \ALIUSPlacedTextContent{90.000}{635.743}{418.395}{ALIUSC1F8135}{\ALIUSFontLatoLight}{12.357}{(Melbourne, Australia). In your research, you discuss contemporary debates of}
  \ALIUSPlacedTextContent{90.000}{653.023}{418.395}{ALIUSC1F8135}{\ALIUSFontLatoLight}{12.357}{philosophy of mind, including consciousness, through the lens of Buddhist and}
  \ALIUSPlacedTextContent{90.000}{670.303}{418.395}{ALIUSC1F8135}{\ALIUSFontLatoLight}{12.357}{Hinduist philosophy. Would you say that your interest for Indian philosophy}
  \ALIUSPlacedTextContent{90.000}{687.583}{418.530}{ALIUSC1F8135}{\ALIUSFontLatoLight}{12.357}{led you to your interest for the philosophy of mind or the reverse? Is there}
  \ALIUSPlacedTextContent{90.000}{704.863}{418.320}{ALIUSC1F8135}{\ALIUSFontLatoLight}{12.357}{something in particular in the Indian philosophical corpus that sparkled your}
  \ALIUSPlacedTextContent{90.000}{722.143}{49.837}{ALIUSC1F8135}{\ALIUSFontLatoLight}{12.357}{interest?}
  \ALIUSPlacedTextContent{90.000}{792.192}{120.165}{ALIUSC767171}{\ALIUSFontCormorantLight}{12.000}{ALIUS Bulletin n°5 (2021)}
  \ALIUSPlacedTextContent{384.252}{792.192}{120.728}{ALIUSC767171}{\ALIUSFontCormorantLight}{12.000}{aliusresearch.org/bulletin}
\end{tikzpicture}
\null
\newpage
% Page 2
\thispagestyle{empty}
\noindent\begin{tikzpicture}[remember picture,overlay,x=1bp,y=1bp,shift={(current page.north west)}]
  \ALIUSPlacedTextContent{90.000}{35.378}{386.261}{ALIUSC767171}{\ALIUSFontCormorantRegular}{13.920}{Monima Chadha – Transformative experiences in Indian philosophy}
  \ALIUSPlacedTextContent{494.640}{35.472}{13.200}{ALIUSC000000}{\ALIUSFontCormorantRegular}{12.000}{30}
  \ALIUSPlacedTextContent{90.000}{72.098}{418.430}{ALIUSC000000}{\ALIUSFontCormorantRegular}{13.920}{I have always been intrigued by issues in philosophy of mind and language,}
  \ALIUSPlacedTextContent{90.000}{91.538}{418.234}{ALIUSC000000}{\ALIUSFontCormorantRegular}{13.920}{so in a sense it was philosophy of mind that came first. My interest in Indian}
  \ALIUSPlacedTextContent{90.000}{110.978}{418.261}{ALIUSC000000}{\ALIUSFontCormorantRegular}{13.920}{philosophy was sparked by the Buddhist no-self doctrine. Some Buddhists,}
  \ALIUSPlacedTextContent{90.000}{130.658}{62.941}{ALIUSC000000}{\ALIUSFontCormorantItalic}{13.920}{Abhidharma}
  \ALIUSPlacedTextContent{152.985}{130.658}{355.236}{ALIUSC000000}{\ALIUSFontCormorantRegular}{13.920}{philosophers in the northern Indian tradition, in particular,}
  \ALIUSPlacedTextContent{90.000}{150.098}{418.243}{ALIUSC000000}{\ALIUSFontCormorantRegular}{13.920}{claim that the work of the self can be transferred to the mind. The self is an}
  \ALIUSPlacedTextContent{90.000}{169.538}{418.267}{ALIUSC000000}{\ALIUSFontCormorantRegular}{13.920}{ontological dangler that does not do any causal work. I became interested in}
  \ALIUSPlacedTextContent{90.000}{188.978}{132.824}{ALIUSC000000}{\ALIUSFontCormorantRegular}{13.920}{investigating this claim.}
  \ALIUSPlacedTextContent{90.000}{220.783}{383.161}{ALIUSC1F8135}{\ALIUSFontLatoLight}{12.357}{An important concern for a lot of traditions in Indian philosophy (both}
  \ALIUSPlacedTextContent{473.776}{220.783}{31.212}{ALIUSC1F8135}{\ALIUSFontLatoLightItalic}{12.357}{Astika}
  \ALIUSPlacedTextContent{90.000}{238.063}{23.224}{ALIUSC1F8135}{\ALIUSFontLatoLight}{12.357}{and}
  \ALIUSPlacedTextContent{113.102}{238.063}{38.621}{ALIUSC1F8135}{\ALIUSFontLatoLightItalic}{12.357}{Nastika}
  \ALIUSPlacedTextContent{151.726}{238.063}{356.675}{ALIUSC1F8135}{\ALIUSFontLatoLight}{12.357}{systems) has been the nature of consciousness.  However, Indian}
  \ALIUSPlacedTextContent{90.000}{255.103}{418.279}{ALIUSC1F8135}{\ALIUSFontLatoLight}{12.357}{philosophy has been totally overlooked in the standard curricula of the}
  \ALIUSPlacedTextContent{90.000}{272.383}{418.332}{ALIUSC1F8135}{\ALIUSFontLatoLight}{12.357}{philosophy of mind despite the fact that a lot of the contemporary discussions}
  \ALIUSPlacedTextContent{90.000}{289.663}{418.486}{ALIUSC1F8135}{\ALIUSFontLatoLight}{12.357}{are about consciousness. In fact, a typical course starts with Descartes and}
  \ALIUSPlacedTextContent{90.000}{306.943}{418.484}{ALIUSC1F8135}{\ALIUSFontLatoLight}{12.357}{ends with discussions about Chalmers' philosophical zombies ignoring more}
  \ALIUSPlacedTextContent{90.000}{324.223}{418.417}{ALIUSC1F8135}{\ALIUSFontLatoLight}{12.357}{than 4000 years of Indian reflections about consciousness. Do you have a}
  \ALIUSPlacedTextContent{90.000}{341.503}{418.388}{ALIUSC1F8135}{\ALIUSFontLatoLight}{12.357}{diagnostic of why this is the case? How could Indian philosophy contribute in}
  \ALIUSPlacedTextContent{90.000}{358.783}{418.504}{ALIUSC1F8135}{\ALIUSFontLatoLight}{12.357}{your opinion to our understanding of the mind and do you have any}
  \ALIUSPlacedTextContent{90.000}{376.063}{418.387}{ALIUSC1F8135}{\ALIUSFontLatoLight}{12.357}{suggestions of Indian texts or readings to implement in the philosophy of mind}
  \ALIUSPlacedTextContent{90.000}{393.103}{168.617}{ALIUSC1F8135}{\ALIUSFontLatoLight}{12.357}{curricula to overcome this gap?}
  \ALIUSPlacedTextContent{90.000}{422.258}{418.258}{ALIUSC000000}{\ALIUSFontCormorantRegular}{13.920}{I hesitate to recommend texts because the Indian tradition is primarily an}
  \ALIUSPlacedTextContent{90.000}{441.938}{418.174}{ALIUSC000000}{\ALIUSFontCormorantRegular}{13.920}{oral tradition. So, the texts are quite cryptic and difficult to follow. For}
  \ALIUSPlacedTextContent{90.000}{461.378}{114.892}{ALIUSC000000}{\ALIUSFontCormorantRegular}{13.920}{example, the famous}
  \ALIUSPlacedTextContent{205.160}{461.378}{33.642}{ALIUSC000000}{\ALIUSFontCormorantItalic}{13.920}{Nyaya}
  \ALIUSPlacedTextContent{238.830}{461.378}{202.294}{ALIUSC000000}{\ALIUSFontCormorantRegular}{13.920}{proof for the existence of the self in}
  \ALIUSPlacedTextContent{441.358}{461.378}{63.605}{ALIUSC000000}{\ALIUSFontCormorantItalic}{13.920}{Nyāya-sūtra}
  \ALIUSPlacedTextContent{90.000}{480.818}{418.247}{ALIUSC000000}{\ALIUSFontCormorantRegular}{13.920}{1.1.10 states that "Desire, aversion, effort, pleasure, pain etc. are the inferential}
  \ALIUSPlacedTextContent{90.000}{500.258}{301.094}{ALIUSC000000}{\ALIUSFontCormorantRegular}{13.920}{signs of the self". Then there is the whole corpus of the}
  \ALIUSPlacedTextContent{391.220}{500.258}{33.642}{ALIUSC000000}{\ALIUSFontCormorantItalic}{13.920}{Nyaya}
  \ALIUSPlacedTextContent{424.890}{500.258}{83.359}{ALIUSC000000}{\ALIUSFontCormorantRegular}{13.920}{commentaries}
  \ALIUSPlacedTextContent{90.000}{519.698}{418.250}{ALIUSC000000}{\ALIUSFontCormorantRegular}{13.920}{on this sutra which itself contain many sophisticated arguments for the}
  \ALIUSPlacedTextContent{90.000}{539.378}{418.239}{ALIUSC000000}{\ALIUSFontCormorantRegular}{13.920}{existence of the self. The other problem is that good translations of the}
  \ALIUSPlacedTextContent{90.000}{558.818}{418.285}{ALIUSC000000}{\ALIUSFontCormorantRegular}{13.920}{primary texts have been few and far between. The situation is changing now.}
  \ALIUSPlacedTextContent{90.000}{578.258}{350.156}{ALIUSC000000}{\ALIUSFontCormorantRegular}{13.920}{Books I would recommend readily for undergraduate curricula.}
  \ALIUSPlacedTextContent{90.000}{609.698}{32.232}{ALIUSC000000}{\ALIUSFontCormorantRegular}{13.920}{1.The}
  \ALIUSPlacedTextContent{123.761}{609.698}{63.605}{ALIUSC000000}{\ALIUSFontCormorantItalic}{13.920}{Nyāya-sūtra}
  \ALIUSPlacedTextContent{187.404}{609.698}{320.834}{ALIUSC000000}{\ALIUSFontCormorantRegular}{13.920}{: Selections with Early Commentaries trans. by Matthew}
  \ALIUSPlacedTextContent{90.000}{629.378}{150.687}{ALIUSC000000}{\ALIUSFontCormorantRegular}{13.920}{Dasti and Stephen Phillips.}
  \ALIUSPlacedTextContent{90.000}{660.818}{296.817}{ALIUSC000000}{\ALIUSFontCormorantRegular}{13.920}{2. The Concealed Art of the Soul, by Jonardon Ganeri}
  \ALIUSPlacedTextContent{90.000}{692.258}{255.257}{ALIUSC000000}{\ALIUSFontCormorantRegular}{13.920}{3. Buddhism as Philosophy by Mark Siderits.}
  \ALIUSPlacedTextContent{90.000}{792.192}{120.165}{ALIUSC767171}{\ALIUSFontCormorantLight}{12.000}{ALIUS Bulletin n°5 (2021)}
  \ALIUSPlacedTextContent{384.252}{792.192}{120.728}{ALIUSC767171}{\ALIUSFontCormorantLight}{12.000}{aliusresearch.org/bulletin}
\end{tikzpicture}
\null
\newpage
% Page 3
\thispagestyle{empty}
\noindent\begin{tikzpicture}[remember picture,overlay,x=1bp,y=1bp,shift={(current page.north west)}]
  \ALIUSPlacedTextContent{90.000}{35.378}{386.261}{ALIUSC767171}{\ALIUSFontCormorantRegular}{13.920}{Monima Chadha – Transformative experiences in Indian philosophy}
  \ALIUSPlacedTextContent{496.320}{35.472}{11.460}{ALIUSC000000}{\ALIUSFontCormorantRegular}{12.000}{31}
  \ALIUSPlacedTextContent{105.600}{92.357}{14.160}{ALIUSC595959}{\ALIUSFontCormorantLight}{48.000}{\ALIUSPullQuoteOpen}
  \ALIUSPlacedTextContent{133.680}{92.357}{324.354}{ALIUSC000000}{\ALIUSFontLatoLight}{15.120}{I hesitate to recommend texts because the Indian}
  \ALIUSPlacedTextContent{468.240}{104.448}{14.160}{ALIUSC595959}{\ALIUSFontCormorantLight}{48.000}{\ALIUSPullQuoteClose}
  \ALIUSPlacedTextContent{133.680}{110.357}{238.018}{ALIUSC000000}{\ALIUSFontLatoLight}{15.120}{tradition is primarily an oral tradition.}
  \ALIUSPlacedTextContent{90.000}{161.743}{418.304}{ALIUSC1F8135}{\ALIUSFontLatoLight}{12.357}{Since the 90's there has been a rising interest in the neuroscience of}
  \ALIUSPlacedTextContent{90.000}{179.023}{418.574}{ALIUSC1F8135}{\ALIUSFontLatoLight}{12.357}{consciousness. In the field, neurobiologists and psychologists usually operate}
  \ALIUSPlacedTextContent{90.000}{196.303}{418.282}{ALIUSC1F8135}{\ALIUSFontLatoLight}{12.357}{with a cognitive concept of consciousness: consciousness is cognitive access}
  \ALIUSPlacedTextContent{90.000}{213.583}{418.373}{ALIUSC1F8135}{\ALIUSFontLatoLight}{12.357}{(roughly: awareness). Some philosophers have objected to this framework that}
  \ALIUSPlacedTextContent{90.000}{230.863}{418.314}{ALIUSC1F8135}{\ALIUSFontLatoLight}{12.357}{a cognitive concept of consciousness isn't enough, to solve the hard problem}
  \ALIUSPlacedTextContent{90.000}{248.143}{418.432}{ALIUSC1F8135}{\ALIUSFontLatoLight}{12.357}{we need to target what's like to be in a certain state (roughly: experience)}
  \ALIUSPlacedTextContent{90.000}{265.183}{418.406}{ALIUSC1F8135}{\ALIUSFontLatoLight}{12.357}{(Chalmers, 2007). Are there different conceptions of consciousness in the}
  \ALIUSPlacedTextContent{90.000}{282.463}{90.353}{ALIUSC1F8135}{\ALIUSFontLatoLight}{12.357}{Indian tradition?}
  \ALIUSPlacedTextContent{90.000}{311.618}{418.340}{ALIUSC000000}{\ALIUSFontCormorantRegular}{13.920}{Our pre-theoretical concept of consciousness is a mongrel concept which}
  \ALIUSPlacedTextContent{90.000}{331.058}{418.187}{ALIUSC000000}{\ALIUSFontCormorantRegular}{13.920}{calls for refinement if it is to do useful work for a theory of consciousness.}
  \ALIUSPlacedTextContent{90.000}{350.738}{418.275}{ALIUSC000000}{\ALIUSFontCormorantRegular}{13.920}{Just like Block and Chalmers in the contemporary scene, the Indian}
  \ALIUSPlacedTextContent{90.000}{370.178}{418.210}{ALIUSC000000}{\ALIUSFontCormorantRegular}{13.920}{philosophers also suggested many refinements and distinctions to be made}
  \ALIUSPlacedTextContent{90.000}{389.618}{418.185}{ALIUSC000000}{\ALIUSFontCormorantRegular}{13.920}{within the concept of consciousness. But I don't think we can map them}
  \ALIUSPlacedTextContent{90.000}{409.058}{418.304}{ALIUSC000000}{\ALIUSFontCormorantRegular}{13.920}{directly onto distinctions drawn in contemporary philosophy. This requires}
  \ALIUSPlacedTextContent{90.000}{428.738}{418.185}{ALIUSC000000}{\ALIUSFontCormorantRegular}{13.920}{careful comparative work which is yet to be done. That said, I do think there}
  \ALIUSPlacedTextContent{90.000}{448.178}{347.593}{ALIUSC000000}{\ALIUSFontCormorantRegular}{13.920}{are echoes of the phenomenal/cognitive distinction in the}
  \ALIUSPlacedTextContent{442.015}{448.178}{62.941}{ALIUSC000000}{\ALIUSFontCormorantItalic}{13.920}{Abhidharma}
  \ALIUSPlacedTextContent{90.000}{467.618}{273.930}{ALIUSC000000}{\ALIUSFontCormorantRegular}{13.920}{tradition in that the notions of feeling (}
  \ALIUSPlacedTextContent{363.888}{467.618}{35.684}{ALIUSC000000}{\ALIUSFontCormorantItalic}{13.920}{vedana}
  \ALIUSPlacedTextContent{399.614}{467.618}{108.643}{ALIUSC000000}{\ALIUSFontCormorantRegular}{13.920}{) and perceptual}
  \ALIUSPlacedTextContent{90.000}{487.058}{88.650}{ALIUSC000000}{\ALIUSFontCormorantRegular}{13.920}{discrimination (}
  \ALIUSPlacedTextContent{178.667}{487.058}{35.670}{ALIUSC000000}{\ALIUSFontCormorantItalic}{13.920}{saṃjñā}
  \ALIUSPlacedTextContent{214.367}{487.058}{293.859}{ALIUSC000000}{\ALIUSFontCormorantRegular}{13.920}{) as components of consciousness, but I do not think}
  \ALIUSPlacedTextContent{90.000}{506.738}{334.461}{ALIUSC000000}{\ALIUSFontCormorantRegular}{13.920}{this clearly maps onto the phenomenal/cognitive distinction.}
  \ALIUSPlacedTextContent{90.000}{538.303}{418.355}{ALIUSC1F8135}{\ALIUSFontLatoLight}{12.357}{A lot of people would agree with the claim that phenomenology makes life}
  \ALIUSPlacedTextContent{90.000}{555.583}{418.389}{ALIUSC1F8135}{\ALIUSFontLatoLight}{12.357}{worth living. A good life is about -among other things- having good}
  \ALIUSPlacedTextContent{90.000}{572.863}{418.443}{ALIUSC1F8135}{\ALIUSFontLatoLight}{12.357}{experiences. Similarly, prolonged experiences of suffering can undermine the}
  \ALIUSPlacedTextContent{90.000}{589.903}{418.487}{ALIUSC1F8135}{\ALIUSFontLatoLight}{12.357}{value of our life. So, there seems to be an important connection between}
  \ALIUSPlacedTextContent{90.000}{607.183}{418.528}{ALIUSC1F8135}{\ALIUSFontLatoLight}{12.357}{consciousness and value. One explanation of this connection is that}
  \ALIUSPlacedTextContent{90.000}{624.463}{418.389}{ALIUSC1F8135}{\ALIUSFontLatoLight}{12.357}{experiences themselves have an intrinsic valenced dimension (experiences}
  \ALIUSPlacedTextContent{90.000}{641.743}{418.336}{ALIUSC1F8135}{\ALIUSFontLatoLight}{12.357}{feel good, bad, pleasant, unpleasant, etc). Furthermore, both consciousness}
  \ALIUSPlacedTextContent{90.000}{659.023}{418.243}{ALIUSC1F8135}{\ALIUSFontLatoLight}{12.357}{and suffering played an important role in some Indian doctrines, for example,}
  \ALIUSPlacedTextContent{90.000}{676.303}{85.182}{ALIUSC1F8135}{\ALIUSFontLatoLight}{12.357}{the doctrine of}
  \ALIUSPlacedTextContent{175.895}{676.303}{42.415}{ALIUSC1F8135}{\ALIUSFontLatoLightItalic}{12.357}{samsara}
  \ALIUSPlacedTextContent{218.312}{676.303}{129.158}{ALIUSC1F8135}{\ALIUSFontLatoLight}{12.357}{in the Upanishads and}
  \ALIUSPlacedTextContent{348.154}{676.303}{31.115}{ALIUSC1F8135}{\ALIUSFontLatoLightItalic}{12.357}{duhka}
  \ALIUSPlacedTextContent{379.290}{676.303}{129.119}{ALIUSC1F8135}{\ALIUSFontLatoLight}{12.357}{in the first noble truth}
  \ALIUSPlacedTextContent{90.000}{693.583}{418.386}{ALIUSC1F8135}{\ALIUSFontLatoLight}{12.357}{of Buddhism. However, with few exceptions (Kriegel, 2019) in contemporary}
  \ALIUSPlacedTextContent{90.000}{710.863}{418.305}{ALIUSC1F8135}{\ALIUSFontLatoLight}{12.357}{philosophy of perception, the study of consciousness is focused on the}
  \ALIUSPlacedTextContent{90.000}{727.903}{333.662}{ALIUSC1F8135}{\ALIUSFontLatoLight}{12.357}{sensory features of experience ignoring its valenced aspects.}
  \ALIUSPlacedTextContent{90.000}{792.192}{120.165}{ALIUSC767171}{\ALIUSFontCormorantLight}{12.000}{ALIUS Bulletin n°5 (2021)}
  \ALIUSPlacedTextContent{384.252}{792.192}{120.728}{ALIUSC767171}{\ALIUSFontCormorantLight}{12.000}{aliusresearch.org/bulletin}
\end{tikzpicture}
\null
\newpage
% Page 4
\thispagestyle{empty}
\noindent\begin{tikzpicture}[remember picture,overlay,x=1bp,y=1bp,shift={(current page.north west)}]
  \ALIUSPlacedTextContent{90.000}{35.378}{386.261}{ALIUSC767171}{\ALIUSFontCormorantRegular}{13.920}{Monima Chadha – Transformative experiences in Indian philosophy}
  \ALIUSPlacedTextContent{495.600}{35.472}{12.300}{ALIUSC000000}{\ALIUSFontCormorantRegular}{12.000}{32}
  \ALIUSPlacedTextContent{90.000}{72.223}{418.383}{ALIUSC1F8135}{\ALIUSFontLatoLight}{12.357}{What do Indian traditions have to say about the relationship between}
  \ALIUSPlacedTextContent{90.000}{89.503}{147.395}{ALIUSC1F8135}{\ALIUSFontLatoLight}{12.357}{consciousness and value?}
  \ALIUSPlacedTextContent{90.000}{118.658}{418.171}{ALIUSC000000}{\ALIUSFontCormorantRegular}{13.920}{A primary concern that motivated most philosophers in ancient India was to}
  \ALIUSPlacedTextContent{90.000}{138.098}{267.885}{ALIUSC000000}{\ALIUSFontCormorantRegular}{13.920}{find the best way forward in an individual person}
  \ALIUSPlacedTextContent{362.574}{138.098}{145.677}{ALIUSC000000}{\ALIUSFontCormorantRegular}{13.920}{s quest for liberation from}
  \ALIUSPlacedTextContent{356.917}{138.557}{9.872}{ALIUSC000000}{\ALIUSFontCormorantRegular}{13.920}{'}
  \ALIUSPlacedTextContent{90.000}{157.538}{418.238}{ALIUSC000000}{\ALIUSFontCormorantRegular}{13.920}{suffering. All living beings are trapped in the cycle of birth and rebirth}
  \ALIUSPlacedTextContent{90.000}{176.978}{4.301}{ALIUSC000000}{\ALIUSFontCormorantRegular}{13.920}{(}
  \ALIUSPlacedTextContent{94.326}{176.978}{40.872}{ALIUSC000000}{\ALIUSFontCormorantItalic}{13.920}{saṃsāra}
  \ALIUSPlacedTextContent{135.233}{176.978}{373.100}{ALIUSC000000}{\ALIUSFontCormorantRegular}{13.920}{) which according to most classical Indian philosophers, is}
  \ALIUSPlacedTextContent{90.000}{196.658}{366.817}{ALIUSC000000}{\ALIUSFontCormorantRegular}{13.920}{characterised by suffering. The highest goal of life is liberation (}
  \ALIUSPlacedTextContent{456.815}{196.658}{31.134}{ALIUSC000000}{\ALIUSFontCormorantItalic}{13.920}{mokṣa}
  \ALIUSPlacedTextContent{487.992}{196.658}{20.259}{ALIUSC000000}{\ALIUSFontCormorantRegular}{13.920}{or}
  \ALIUSPlacedTextContent{90.000}{216.098}{39.410}{ALIUSC000000}{\ALIUSFontCormorantItalic}{13.920}{nirvāṇa}
  \ALIUSPlacedTextContent{129.437}{216.098}{72.167}{ALIUSC000000}{\ALIUSFontCormorantRegular}{13.920}{). Except the}
  \ALIUSPlacedTextContent{202.047}{216.098}{43.032}{ALIUSC000000}{\ALIUSFontCormorantItalic}{13.920}{Cārvāka}
  \ALIUSPlacedTextContent{245.124}{216.098}{263.007}{ALIUSC000000}{\ALIUSFontCormorantRegular}{13.920}{materialists who believed that death is the end}
  \ALIUSPlacedTextContent{90.000}{235.538}{418.233}{ALIUSC000000}{\ALIUSFontCormorantRegular}{13.920}{and there is nothing like rebirth or liberation, all other philosophical schools}
  \ALIUSPlacedTextContent{90.000}{254.978}{418.252}{ALIUSC000000}{\ALIUSFontCormorantRegular}{13.920}{believed in the possibility of liberation in this life or in future lifetimes. The}
  \ALIUSPlacedTextContent{90.000}{274.658}{418.288}{ALIUSC000000}{\ALIUSFontCormorantRegular}{13.920}{ultimate aim of the classical Indian philosophy teachings is thus to help}
  \ALIUSPlacedTextContent{90.000}{294.098}{418.263}{ALIUSC000000}{\ALIUSFontCormorantRegular}{13.920}{individual persons attain liberation or at least a better life in this life and}
  \ALIUSPlacedTextContent{90.000}{313.538}{418.193}{ALIUSC000000}{\ALIUSFontCormorantRegular}{13.920}{future lifetimes. Most classical Indian philosophers agreed that our ignorance}
  \ALIUSPlacedTextContent{90.000}{332.978}{418.129}{ALIUSC000000}{\ALIUSFontCormorantRegular}{13.920}{about "who we really are" is the source and the means to bringing an end to}
  \ALIUSPlacedTextContent{90.000}{352.658}{418.260}{ALIUSC000000}{\ALIUSFontCormorantRegular}{13.920}{suffering. Thus, metaphysical debates about the nature of consciousness and}
  \ALIUSPlacedTextContent{90.000}{372.098}{418.396}{ALIUSC000000}{\ALIUSFontCormorantRegular}{13.920}{the universe and our place in it are central to the classical Indian}
  \ALIUSPlacedTextContent{90.000}{391.538}{418.257}{ALIUSC000000}{\ALIUSFontCormorantRegular}{13.920}{philosophical traditions, but only insofar as they suggest a route to liberation.}
  \ALIUSPlacedTextContent{90.000}{423.103}{418.457}{ALIUSC1F8135}{\ALIUSFontLatoLight}{12.357}{Current debates on consciousness tackle the possibility that consciousness is}
  \ALIUSPlacedTextContent{90.000}{440.383}{418.376}{ALIUSC1F8135}{\ALIUSFontLatoLight}{12.357}{an illusion (Frankish, 2016). The notion of illusion has been deeply anchored}
  \ALIUSPlacedTextContent{90.000}{457.663}{418.458}{ALIUSC1F8135}{\ALIUSFontLatoLight}{12.357}{in the Indian philosophical tradition, such as the concept of maya, the idea that}
  \ALIUSPlacedTextContent{90.000}{474.943}{336.024}{ALIUSC1F8135}{\ALIUSFontLatoLight}{12.357}{the phenomenal world is illusory, being central for the Hinduist}
  \ALIUSPlacedTextContent{425.330}{474.943}{39.250}{ALIUSC1F8135}{\ALIUSFontLatoLightItalic}{12.357}{Advaita}
  \ALIUSPlacedTextContent{464.591}{474.943}{43.793}{ALIUSC1F8135}{\ALIUSFontLatoLight}{12.357}{school.}
  \ALIUSPlacedTextContent{90.000}{492.223}{418.400}{ALIUSC1F8135}{\ALIUSFontLatoLight}{12.357}{Do you see a way by which discussions in Indian philosophy on consciousness}
  \ALIUSPlacedTextContent{90.000}{509.503}{418.564}{ALIUSC1F8135}{\ALIUSFontLatoLight}{12.357}{echo and differ from contemporary debates about the nature of}
  \ALIUSPlacedTextContent{90.000}{526.783}{85.787}{ALIUSC1F8135}{\ALIUSFontLatoLight}{12.357}{consciousness?}
  \ALIUSPlacedTextContent{90.000}{555.698}{334.574}{ALIUSC000000}{\ALIUSFontCormorantRegular}{13.920}{You are right the notion of illusion gets much airtime in the}
  \ALIUSPlacedTextContent{425.154}{555.698}{40.360}{ALIUSC000000}{\ALIUSFontCormorantItalic}{13.920}{Advaita}
  \ALIUSPlacedTextContent{465.557}{555.698}{42.691}{ALIUSC000000}{\ALIUSFontCormorantRegular}{13.920}{Hindu}
  \ALIUSPlacedTextContent{90.000}{575.378}{418.299}{ALIUSC000000}{\ALIUSFontCormorantRegular}{13.920}{tradition, but it has an equally important role in all other Indian traditions}
  \ALIUSPlacedTextContent{90.000}{594.818}{418.279}{ALIUSC000000}{\ALIUSFontCormorantRegular}{13.920}{as well. This should come as no surprise, philosophers in all traditions have}
  \ALIUSPlacedTextContent{90.000}{614.258}{418.246}{ALIUSC000000}{\ALIUSFontCormorantRegular}{13.920}{been interested in illusions in the discussions of perception, consciousness}
  \ALIUSPlacedTextContent{90.000}{633.698}{418.275}{ALIUSC000000}{\ALIUSFontCormorantRegular}{13.920}{and ontology. The discussions in Indian philosophy differ because the}
  \ALIUSPlacedTextContent{90.000}{653.378}{418.332}{ALIUSC000000}{\ALIUSFontCormorantRegular}{13.920}{primary aim guiding Indian philosophers is how to achieve enlightenment or}
  \ALIUSPlacedTextContent{110.160}{692.596}{14.160}{ALIUSC595959}{\ALIUSFontCormorantLight}{48.000}{\ALIUSPullQuoteOpen}
  \ALIUSPlacedTextContent{138.000}{692.596}{326.575}{ALIUSC000000}{\ALIUSFontLatoLight}{15.120}{The discussions in Indian philosophy differ}
  \ALIUSPlacedTextContent{138.000}{710.596}{57.157}{ALIUSC000000}{\ALIUSFontLatoLight}{15.120}{because}
  \ALIUSPlacedTextContent{209.046}{710.596}{25.474}{ALIUSC000000}{\ALIUSFontLatoLight}{15.120}{the}
  \ALIUSPlacedTextContent{248.409}{710.596}{53.331}{ALIUSC000000}{\ALIUSFontLatoLight}{15.120}{primary}
  \ALIUSPlacedTextContent{315.630}{710.596}{26.684}{ALIUSC000000}{\ALIUSFontLatoLight}{15.120}{aim}
  \ALIUSPlacedTextContent{356.445}{710.596}{50.912}{ALIUSC000000}{\ALIUSFontLatoLight}{15.120}{guiding}
  \ALIUSPlacedTextContent{421.488}{710.596}{43.237}{ALIUSC000000}{\ALIUSFontLatoLight}{15.120}{Indian}
  \ALIUSPlacedTextContent{138.000}{728.596}{326.746}{ALIUSC000000}{\ALIUSFontLatoLight}{15.120}{philosophers is how to achieve enlightenment or}
  \ALIUSPlacedTextContent{477.360}{737.088}{14.160}{ALIUSC595959}{\ALIUSFontCormorantLight}{48.000}{\ALIUSPullQuoteClose}
  \ALIUSPlacedTextContent{138.000}{746.596}{64.018}{ALIUSC000000}{\ALIUSFontLatoLight}{15.120}{liberation.}
  \ALIUSPlacedTextContent{90.000}{792.192}{120.165}{ALIUSC767171}{\ALIUSFontCormorantLight}{12.000}{ALIUS Bulletin n°5 (2021)}
  \ALIUSPlacedTextContent{384.252}{792.192}{120.728}{ALIUSC767171}{\ALIUSFontCormorantLight}{12.000}{aliusresearch.org/bulletin}
\end{tikzpicture}
\null
\newpage
% Page 5
\thispagestyle{empty}
\noindent\begin{tikzpicture}[remember picture,overlay,x=1bp,y=1bp,shift={(current page.north west)}]
  \ALIUSPlacedTextContent{90.000}{35.378}{386.261}{ALIUSC767171}{\ALIUSFontCormorantRegular}{13.920}{Monima Chadha – Transformative experiences in Indian philosophy}
  \ALIUSPlacedTextContent{495.600}{35.472}{12.144}{ALIUSC000000}{\ALIUSFontCormorantRegular}{12.000}{33}
  \ALIUSPlacedTextContent{90.000}{72.098}{418.275}{ALIUSC000000}{\ALIUSFontCormorantRegular}{13.920}{liberation. This concern gives the Indian discussion of consciousness a}
  \ALIUSPlacedTextContent{90.000}{91.538}{418.269}{ALIUSC000000}{\ALIUSFontCormorantRegular}{13.920}{different flavour, but in the end, this concern cannot be divorced from the}
  \ALIUSPlacedTextContent{90.000}{110.978}{418.185}{ALIUSC000000}{\ALIUSFontCormorantRegular}{13.920}{nature of human beings and what they are capable of. So, the Indian}
  \ALIUSPlacedTextContent{90.000}{130.658}{418.203}{ALIUSC000000}{\ALIUSFontCormorantRegular}{13.920}{philosophers are indirectly led to metaphysical and epistemological questions}
  \ALIUSPlacedTextContent{90.000}{150.098}{418.185}{ALIUSC000000}{\ALIUSFontCormorantRegular}{13.920}{about the nature of the world and human beings within it. Keeping this}
  \ALIUSPlacedTextContent{90.000}{169.538}{418.210}{ALIUSC000000}{\ALIUSFontCormorantRegular}{13.920}{primary aim in mind, we should consider what the Indian philosophers have}
  \ALIUSPlacedTextContent{90.000}{188.978}{418.254}{ALIUSC000000}{\ALIUSFontCormorantRegular}{13.920}{to say about the mind and we will see many of the same questions being}
  \ALIUSPlacedTextContent{90.000}{208.658}{222.395}{ALIUSC000000}{\ALIUSFontCormorantRegular}{13.920}{discussed in the classical Indian debates.}
  \ALIUSPlacedTextContent{90.000}{240.223}{373.018}{ALIUSC1F8135}{\ALIUSFontLatoLight}{12.357}{One of the most important doctrines in the Upanishads and the}
  \ALIUSPlacedTextContent{465.739}{240.223}{39.250}{ALIUSC1F8135}{\ALIUSFontLatoLightItalic}{12.357}{Advaita}
  \ALIUSPlacedTextContent{90.000}{257.503}{43.081}{ALIUSC1F8135}{\ALIUSFontLatoLightItalic}{12.357}{Vedanta}
  \ALIUSPlacedTextContent{133.099}{257.503}{136.553}{ALIUSC1F8135}{\ALIUSFontLatoLight}{12.357}{is the identity between}
  \ALIUSPlacedTextContent{269.651}{257.503}{51.910}{ALIUSC1F8135}{\ALIUSFontLatoLightItalic}{12.357}{Brahman}
  \ALIUSPlacedTextContent{321.598}{257.503}{28.932}{ALIUSC1F8135}{\ALIUSFontLatoLight}{12.357}{and}
  \ALIUSPlacedTextContent{352.843}{257.503}{34.320}{ALIUSC1F8135}{\ALIUSFontLatoLightItalic}{12.357}{Atman}
  \ALIUSPlacedTextContent{387.174}{257.503}{121.184}{ALIUSC1F8135}{\ALIUSFontLatoLight}{12.357}{(self). A tendentious}
  \ALIUSPlacedTextContent{90.000}{274.783}{418.425}{ALIUSC1F8135}{\ALIUSFontLatoLight}{12.357}{reading of this claim is that fundamental reality involves consciousness. Similar}
  \ALIUSPlacedTextContent{90.000}{292.063}{280.136}{ALIUSC1F8135}{\ALIUSFontLatoLight}{12.357}{claims are found in unorthodox systems, like}
  \ALIUSPlacedTextContent{376.337}{292.063}{62.400}{ALIUSC1F8135}{\ALIUSFontLatoLightItalic}{12.357}{Abhidharma}
  \ALIUSPlacedTextContent{438.748}{292.063}{69.651}{ALIUSC1F8135}{\ALIUSFontLatoLight}{12.357}{Buddhism.}
  \ALIUSPlacedTextContent{90.000}{309.103}{418.302}{ALIUSC1F8135}{\ALIUSFontLatoLight}{12.357}{Furthermore, in the contemporary metaphysics of mind, Russellian Monists}
  \ALIUSPlacedTextContent{90.000}{326.383}{418.310}{ALIUSC1F8135}{\ALIUSFontLatoLight}{12.357}{have championed the idea that fundamental reality involves consciousness (or}
  \ALIUSPlacedTextContent{90.000}{343.663}{418.330}{ALIUSC1F8135}{\ALIUSFontLatoLight}{12.357}{proto-phenomenal properties) in a widespread manner (Goff, 2017). What are}
  \ALIUSPlacedTextContent{90.000}{360.943}{180.785}{ALIUSC1F8135}{\ALIUSFontLatoLight}{12.357}{the main differences between}
  \ALIUSPlacedTextContent{276.199}{360.943}{43.081}{ALIUSC1F8135}{\ALIUSFontLatoLightItalic}{12.357}{Vedanta}
  \ALIUSPlacedTextContent{319.297}{360.943}{117.840}{ALIUSC1F8135}{\ALIUSFontLatoLight}{12.357}{'s panpsychism and}
  \ALIUSPlacedTextContent{442.589}{360.943}{62.400}{ALIUSC1F8135}{\ALIUSFontLatoLightItalic}{12.357}{Abhidharma}
  \ALIUSPlacedTextContent{90.000}{378.223}{418.393}{ALIUSC1F8135}{\ALIUSFontLatoLight}{12.357}{panpsychism? Do any of these views have something to say about the}
  \ALIUSPlacedTextContent{90.000}{395.503}{294.126}{ALIUSC1F8135}{\ALIUSFontLatoLight}{12.357}{combination problem for contemporary panpsychism?}
  \ALIUSPlacedTextContent{90.000}{424.658}{418.221}{ALIUSC000000}{\ALIUSFontCormorantRegular}{13.920}{Vedanta panpsychism is rooted in their monism: consciousness is the only}
  \ALIUSPlacedTextContent{90.000}{444.098}{40.922}{ALIUSC000000}{\ALIUSFontCormorantRegular}{13.920}{reality.}
  \ALIUSPlacedTextContent{131.696}{444.098}{62.941}{ALIUSC000000}{\ALIUSFontCormorantItalic}{13.920}{Abhidharma}
  \ALIUSPlacedTextContent{194.680}{444.098}{313.598}{ALIUSC000000}{\ALIUSFontCormorantRegular}{13.920}{Buddhism accepts physical and mental atoms as part of}
  \ALIUSPlacedTextContent{90.000}{463.538}{418.277}{ALIUSC000000}{\ALIUSFontCormorantRegular}{13.920}{the basic furniture of the universe. This, I think, is the most important}
  \ALIUSPlacedTextContent{90.000}{482.978}{418.151}{ALIUSC000000}{\ALIUSFontCormorantRegular}{13.920}{difference. I am not sure what Vedantins have to say about the combination}
  \ALIUSPlacedTextContent{90.000}{502.658}{291.695}{ALIUSC000000}{\ALIUSFontCormorantRegular}{13.920}{problem, but there is a new special volume of the}
  \ALIUSPlacedTextContent{384.110}{502.658}{34.104}{ALIUSC000000}{\ALIUSFontCormorantItalic}{13.920}{Monist}
  \ALIUSPlacedTextContent{418.255}{502.658}{89.996}{ALIUSC000000}{\ALIUSFontCormorantRegular}{13.920}{coming out in}
  \ALIUSPlacedTextContent{90.000}{522.098}{418.245}{ALIUSC000000}{\ALIUSFontCormorantRegular}{13.920}{January 2022 which addresses Cosmopsychism (the holistic counterpart of}
  \ALIUSPlacedTextContent{90.000}{541.538}{418.336}{ALIUSC000000}{\ALIUSFontCormorantRegular}{13.920}{panpsychism) and Indian philosophy might have more thoughtful responses}
  \ALIUSPlacedTextContent{90.000}{560.978}{418.253}{ALIUSC000000}{\ALIUSFontCormorantRegular}{13.920}{to this issue. I have a paper forthcoming in that volume which suggests that}
  \ALIUSPlacedTextContent{90.000}{580.658}{20.681}{ALIUSC000000}{\ALIUSFontCormorantRegular}{13.920}{the}
  \ALIUSPlacedTextContent{112.795}{580.658}{62.941}{ALIUSC000000}{\ALIUSFontCormorantItalic}{13.920}{Abhidharma}
  \ALIUSPlacedTextContent{175.780}{580.658}{332.505}{ALIUSC000000}{\ALIUSFontCormorantRegular}{13.920}{philosophers have various resources that they may use to}
  \ALIUSPlacedTextContent{90.000}{600.098}{418.283}{ALIUSC000000}{\ALIUSFontCormorantRegular}{13.920}{respond to the group of problems that are called combination problems: the}
  \ALIUSPlacedTextContent{90.000}{619.538}{376.825}{ALIUSC000000}{\ALIUSFontCormorantRegular}{13.920}{subject combination problem, the quality combination problem, etc.}
  \ALIUSPlacedTextContent{90.000}{651.103}{158.092}{ALIUSC1F8135}{\ALIUSFontLatoLight}{12.357}{You recently explained how}
  \ALIUSPlacedTextContent{249.757}{651.103}{62.405}{ALIUSC1F8135}{\ALIUSFontLatoLightItalic}{12.357}{Abhidharma}
  \ALIUSPlacedTextContent{312.168}{651.103}{196.203}{ALIUSC1F8135}{\ALIUSFontLatoLight}{12.357}{Buddhism had a panprotopsychist}
  \ALIUSPlacedTextContent{90.000}{668.383}{232.116}{ALIUSC1F8135}{\ALIUSFontLatoLight}{12.357}{account of reality based on the notion of}
  \ALIUSPlacedTextContent{323.347}{668.383}{42.367}{ALIUSC1F8135}{\ALIUSFontLatoLightItalic}{12.357}{dharma}
  \ALIUSPlacedTextContent{366.868}{668.383}{90.376}{ALIUSC1F8135}{\ALIUSFontLatoLight}{12.357}{(Chadha, 2019).}
  \ALIUSPlacedTextContent{458.451}{668.383}{49.704}{ALIUSC1F8135}{\ALIUSFontLatoLightItalic}{12.357}{Dharmas}
  \ALIUSPlacedTextContent{90.000}{685.663}{119.504}{ALIUSC1F8135}{\ALIUSFontLatoLight}{12.357}{can be described as}
  \ALIUSPlacedTextContent{217.020}{685.663}{291.383}{ALIUSC1F8135}{\ALIUSFontLatoLight}{12.357}{particular qualities of existence and they can be}
  \ALIUSPlacedTextContent{90.000}{702.943}{418.369}{ALIUSC1F8135}{\ALIUSFontLatoLight}{12.357}{aggregated into subjective experiences under the combined action of mental}
  \ALIUSPlacedTextContent{90.000}{720.223}{418.399}{ALIUSC1F8135}{\ALIUSFontLatoLight}{12.357}{processes. Our attention has been brought on the suggestion that differences}
  \ALIUSPlacedTextContent{90.000}{737.503}{418.399}{ALIUSC1F8135}{\ALIUSFontLatoLight}{12.357}{in metaphysical conception about consciousness may depend on practices, as}
  \ALIUSPlacedTextContent{90.000}{754.783}{47.581}{ALIUSC1F8135}{\ALIUSFontLatoLight}{12.357}{you say:}
  \ALIUSPlacedTextContent{90.000}{792.192}{120.165}{ALIUSC767171}{\ALIUSFontCormorantLight}{12.000}{ALIUS Bulletin n°5 (2021)}
  \ALIUSPlacedTextContent{384.252}{792.192}{120.728}{ALIUSC767171}{\ALIUSFontCormorantLight}{12.000}{aliusresearch.org/bulletin}
\end{tikzpicture}
\null
\newpage
% Page 6
\thispagestyle{empty}
\noindent\begin{tikzpicture}[remember picture,overlay,x=1bp,y=1bp,shift={(current page.north west)}]
  \ALIUSPlacedTextContent{90.000}{35.378}{386.261}{ALIUSC767171}{\ALIUSFontCormorantRegular}{13.920}{Monima Chadha – Transformative experiences in Indian philosophy}
  \ALIUSPlacedTextContent{494.880}{35.472}{12.888}{ALIUSC000000}{\ALIUSFontCormorantRegular}{12.000}{34}
  \ALIUSPlacedTextContent{90.000}{72.223}{418.423}{ALIUSC1F8135}{\ALIUSFontLatoLight}{12.357}{"If we think that conscious states supervene on collections of present mental}
  \ALIUSPlacedTextContent{90.000}{89.503}{47.288}{ALIUSC1F8135}{\ALIUSFontLatoLightItalic}{12.357}{dharmas}
  \ALIUSPlacedTextContent{138.712}{89.503}{369.684}{ALIUSC1F8135}{\ALIUSFontLatoLight}{12.357}{which are best thought of as proto-conscious or proto-intentional}
  \ALIUSPlacedTextContent{90.000}{106.783}{418.342}{ALIUSC1F8135}{\ALIUSFontLatoLight}{12.357}{features as the result of logical analysis, then we favor the panprotopsychist}
  \ALIUSPlacedTextContent{90.000}{124.063}{237.045}{ALIUSC1F8135}{\ALIUSFontLatoLight}{12.357}{option. Alternatively, if we hold that mental}
  \ALIUSPlacedTextContent{327.107}{124.063}{50.469}{ALIUSC1F8135}{\ALIUSFontLatoLightItalic}{12.357}{dharmas}
  \ALIUSPlacedTextContent{378.000}{124.063}{126.986}{ALIUSC1F8135}{\ALIUSFontLatoLight}{12.357}{are potentially pheno-}
  \ALIUSPlacedTextContent{90.000}{141.343}{418.409}{ALIUSC1F8135}{\ALIUSFontLatoLight}{12.357}{menologically available as they can be discerned as such by experts who have}
  \ALIUSPlacedTextContent{90.000}{158.383}{418.375}{ALIUSC1F8135}{\ALIUSFontLatoLight}{12.357}{mastered the art of mindfulness  meditation, then we favour the panpsychist}
  \ALIUSPlacedTextContent{90.000}{175.663}{85.911}{ALIUSC1F8135}{\ALIUSFontLatoLight}{12.357}{option." (p.31).}
  \ALIUSPlacedTextContent{90.000}{192.943}{418.396}{ALIUSC1F8135}{\ALIUSFontLatoLight}{12.357}{Could you elaborate on how Buddhists have conceptualized the interrelation}
  \ALIUSPlacedTextContent{90.000}{210.223}{418.395}{ALIUSC1F8135}{\ALIUSFontLatoLight}{12.357}{between understanding consciousness and experiencing conscious states?}
  \ALIUSPlacedTextContent{90.000}{227.503}{418.428}{ALIUSC1F8135}{\ALIUSFontLatoLight}{12.357}{How can the experience of conscious states such as meditation change our}
  \ALIUSPlacedTextContent{90.000}{244.783}{228.846}{ALIUSC1F8135}{\ALIUSFontLatoLight}{12.357}{ideas about the nature of consciousness?}
  \ALIUSPlacedTextContent{90.000}{273.938}{418.264}{ALIUSC000000}{\ALIUSFontCormorantRegular}{13.920}{This is a difficult question to answer because there is not one Buddhist view.}
  \ALIUSPlacedTextContent{90.000}{293.378}{418.226}{ALIUSC000000}{\ALIUSFontCormorantRegular}{13.920}{There are many different Buddhist traditions, and they have different views}
  \ALIUSPlacedTextContent{90.000}{312.818}{418.070}{ALIUSC000000}{\ALIUSFontCormorantRegular}{13.920}{about conscious states. I think that it is worth emphasising is that there is}
  \ALIUSPlacedTextContent{90.000}{332.258}{418.243}{ALIUSC000000}{\ALIUSFontCormorantRegular}{13.920}{nothing like mind or consciousness, all there is, is a collection of mental states}
  \ALIUSPlacedTextContent{90.000}{351.938}{418.215}{ALIUSC000000}{\ALIUSFontCormorantRegular}{13.920}{and conscious experiences. And again, there is not one kind of meditation}
  \ALIUSPlacedTextContent{90.000}{371.378}{418.194}{ALIUSC000000}{\ALIUSFontCormorantRegular}{13.920}{but different meditation practices depending on the tradition. To take an}
  \ALIUSPlacedTextContent{90.000}{390.818}{282.650}{ALIUSC000000}{\ALIUSFontCormorantRegular}{13.920}{example, in Vasubandhu's seminal text the}
  \ALIUSPlacedTextContent{382.306}{390.818}{119.526}{ALIUSC000000}{\ALIUSFontCormorantItalic}{13.920}{Abhidharmakośabhasya}
  \ALIUSPlacedTextContent{501.892}{390.818}{6.366}{ALIUSC000000}{\ALIUSFontCormorantRegular}{13.920}{,}
  \ALIUSPlacedTextContent{90.000}{410.258}{418.259}{ALIUSC000000}{\ALIUSFontCormorantRegular}{13.920}{mindfulness meditation is explained carefully as a stepwise progression.}
  \ALIUSPlacedTextContent{90.000}{429.938}{418.210}{ALIUSC000000}{\ALIUSFontCormorantRegular}{13.920}{Mindfulness meditation has the aim of curbing thoughts that proliferate}
  \ALIUSPlacedTextContent{90.000}{449.378}{418.240}{ALIUSC000000}{\ALIUSFontCormorantRegular}{13.920}{naturally because of the variety of external objects and karmic imprints. This}
  \ALIUSPlacedTextContent{90.000}{468.818}{418.254}{ALIUSC000000}{\ALIUSFontCormorantRegular}{13.920}{aim is achieved by controlling the mind by focusing attention on breathing}
  \ALIUSPlacedTextContent{90.000}{488.258}{418.269}{ALIUSC000000}{\ALIUSFontCormorantRegular}{13.920}{to eliminate mind-wandering. The initial aim of mindfulness meditation is}
  \ALIUSPlacedTextContent{90.000}{507.698}{418.347}{ALIUSC000000}{\ALIUSFontCormorantRegular}{13.920}{to train the mind to fix attention on the complex phenomenon of breathing}
  \ALIUSPlacedTextContent{90.000}{527.378}{418.210}{ALIUSC000000}{\ALIUSFontCormorantRegular}{13.920}{in order to analyze or break down the phenomenon into its most basic}
  \ALIUSPlacedTextContent{90.000}{546.818}{77.189}{ALIUSC000000}{\ALIUSFontCormorantRegular}{13.920}{constituents (}
  \ALIUSPlacedTextContent{167.254}{546.818}{43.013}{ALIUSC000000}{\ALIUSFontCormorantItalic}{13.920}{dharmas}
  \ALIUSPlacedTextContent{210.303}{546.818}{297.880}{ALIUSC000000}{\ALIUSFontCormorantRegular}{13.920}{), so that the meditator discerns the real nature of}
  \ALIUSPlacedTextContent{90.000}{566.258}{418.226}{ALIUSC000000}{\ALIUSFontCormorantRegular}{13.920}{conscious experiences. The process of meditation, however, reveals that the}
  \ALIUSPlacedTextContent{90.000}{585.698}{225.390}{ALIUSC000000}{\ALIUSFontCormorantRegular}{13.920}{realization of the real nature of ultimate}
  \ALIUSPlacedTextContent{315.409}{585.698}{43.013}{ALIUSC000000}{\ALIUSFontCormorantItalic}{13.920}{dharmas}
  \ALIUSPlacedTextContent{358.458}{585.698}{149.770}{ALIUSC000000}{\ALIUSFontCormorantRegular}{13.920}{(impermanence, suffering,}
  \ALIUSPlacedTextContent{90.000}{605.378}{418.246}{ALIUSC000000}{\ALIUSFontCormorantRegular}{13.920}{emptiness, and being not-self) results in an attitudinal and behavioural}
  \ALIUSPlacedTextContent{90.000}{624.818}{418.296}{ALIUSC000000}{\ALIUSFontCormorantRegular}{13.920}{change. The knowledge of real nature of self and other objects leads to the}
  \ALIUSPlacedTextContent{90.000}{644.258}{291.257}{ALIUSC000000}{\ALIUSFontCormorantRegular}{13.920}{cessation of all emotional attachments and desires.}
  \ALIUSPlacedTextContent{90.000}{683.503}{418.379}{ALIUSC1F8135}{\ALIUSFontLatoLight}{12.357}{Both neuroscientists and philosophers recognize the role of introspection and}
  \ALIUSPlacedTextContent{90.000}{700.783}{418.369}{ALIUSC1F8135}{\ALIUSFontLatoLight}{12.357}{first-person data for a science of consciousness. What's the role of}
  \ALIUSPlacedTextContent{90.000}{717.823}{418.240}{ALIUSC1F8135}{\ALIUSFontLatoLight}{12.357}{introspection and first-person data to understand consciousness in the Indian}
  \ALIUSPlacedTextContent{90.000}{735.103}{418.423}{ALIUSC1F8135}{\ALIUSFontLatoLight}{12.357}{traditions? Does the Indian philosophical legacy have something to contribute}
  \ALIUSPlacedTextContent{90.000}{752.383}{344.166}{ALIUSC1F8135}{\ALIUSFontLatoLight}{12.357}{to the prospects of a scientific understanding of consciousness?}
  \ALIUSPlacedTextContent{90.000}{792.192}{120.165}{ALIUSC767171}{\ALIUSFontCormorantLight}{12.000}{ALIUS Bulletin n°5 (2021)}
  \ALIUSPlacedTextContent{384.252}{792.192}{120.728}{ALIUSC767171}{\ALIUSFontCormorantLight}{12.000}{aliusresearch.org/bulletin}
\end{tikzpicture}
\null
\newpage
% Page 7
\thispagestyle{empty}
\noindent\begin{tikzpicture}[remember picture,overlay,x=1bp,y=1bp,shift={(current page.north west)}]
  \ALIUSPlacedTextContent{90.000}{35.378}{386.261}{ALIUSC767171}{\ALIUSFontCormorantRegular}{13.920}{Monima Chadha – Transformative experiences in Indian philosophy}
  \ALIUSPlacedTextContent{495.360}{35.472}{12.384}{ALIUSC000000}{\ALIUSFontCormorantRegular}{12.000}{35}
  \ALIUSPlacedTextContent{115.200}{88.517}{14.160}{ALIUSC595959}{\ALIUSFontCormorantLight}{48.000}{\ALIUSPullQuoteOpen}
  \ALIUSPlacedTextContent{141.840}{88.517}{317.611}{ALIUSC000000}{\ALIUSFontLatoLight}{15.120}{Indian philosophers have an additional resource}
  \ALIUSPlacedTextContent{141.840}{106.517}{317.638}{ALIUSC000000}{\ALIUSFontLatoLight}{15.120}{in first-person data available in meditative}
  \ALIUSPlacedTextContent{463.680}{112.848}{14.160}{ALIUSC595959}{\ALIUSFontCormorantLight}{48.000}{\ALIUSPullQuoteClose}
  \ALIUSPlacedTextContent{141.840}{124.517}{80.578}{ALIUSC000000}{\ALIUSFontLatoLight}{15.120}{experiences.}
  \ALIUSPlacedTextContent{90.000}{171.938}{418.258}{ALIUSC000000}{\ALIUSFontCormorantRegular}{13.920}{Philosophers cannot afford to ignore first-person data and introspection if}
  \ALIUSPlacedTextContent{90.000}{191.378}{418.250}{ALIUSC000000}{\ALIUSFontCormorantRegular}{13.920}{they want to understand the nature of consciousness. Indian philosophers}
  \ALIUSPlacedTextContent{90.000}{210.818}{418.248}{ALIUSC000000}{\ALIUSFontCormorantRegular}{13.920}{have an additional resource in first-person data available in meditative}
  \ALIUSPlacedTextContent{90.000}{230.258}{418.278}{ALIUSC000000}{\ALIUSFontCormorantRegular}{13.920}{experiences. The "special access" to one's own conscious states available in}
  \ALIUSPlacedTextContent{90.000}{249.938}{418.249}{ALIUSC000000}{\ALIUSFontCormorantRegular}{13.920}{conscious experiences is amenable to be investigated by third-person}
  \ALIUSPlacedTextContent{90.000}{269.378}{418.346}{ALIUSC000000}{\ALIUSFontCormorantRegular}{13.920}{methods in neuroscience of consciousness. This has the potential to offer}
  \ALIUSPlacedTextContent{90.000}{288.818}{418.193}{ALIUSC000000}{\ALIUSFontCormorantRegular}{13.920}{additional insights that might contribute to a scientific understanding of}
  \ALIUSPlacedTextContent{90.000}{308.258}{418.199}{ALIUSC000000}{\ALIUSFontCormorantRegular}{13.920}{consciousness. But it needs a deeper understanding of the particular Indian}
  \ALIUSPlacedTextContent{90.000}{327.938}{418.355}{ALIUSC000000}{\ALIUSFontCormorantRegular}{13.920}{tradition in question and the role that meditation practices are supposed to}
  \ALIUSPlacedTextContent{90.000}{347.378}{225.692}{ALIUSC000000}{\ALIUSFontCormorantRegular}{13.920}{lead to intellectual and spiritual progress.}
  \ALIUSPlacedTextContent{90.000}{378.943}{176.577}{ALIUSC1F8135}{\ALIUSFontLatoLight}{12.357}{The practice of meditation (}
  \ALIUSPlacedTextContent{266.632}{378.943}{37.621}{ALIUSC1F8135}{\ALIUSFontLatoLightItalic}{12.357}{dhyana}
  \ALIUSPlacedTextContent{304.281}{378.943}{204.059}{ALIUSC1F8135}{\ALIUSFontLatoLight}{12.357}{) is pervasive among the Indian}
  \ALIUSPlacedTextContent{90.000}{396.223}{418.376}{ALIUSC1F8135}{\ALIUSFontLatoLight}{12.357}{philosophical tradition. Through meditation, altered states of consciousness}
  \ALIUSPlacedTextContent{90.000}{413.503}{418.337}{ALIUSC1F8135}{\ALIUSFontLatoLight}{12.357}{can be attained through which fundamental insight about the nature of reality}
  \ALIUSPlacedTextContent{90.000}{430.783}{288.623}{ALIUSC1F8135}{\ALIUSFontLatoLight}{12.357}{can be attained. In some traditions, this is known as '}
  \ALIUSPlacedTextContent{378.600}{430.783}{42.714}{ALIUSC1F8135}{\ALIUSFontLatoLightItalic}{12.357}{samadhi}
  \ALIUSPlacedTextContent{421.306}{430.783}{87.093}{ALIUSC1F8135}{\ALIUSFontLatoLight}{12.357}{'. The very idea}
  \ALIUSPlacedTextContent{90.000}{448.063}{14.637}{ALIUSC1F8135}{\ALIUSFontLatoLight}{12.357}{of}
  \ALIUSPlacedTextContent{106.369}{448.063}{37.924}{ALIUSC1F8135}{\ALIUSFontLatoLightItalic}{12.357}{nirvana}
  \ALIUSPlacedTextContent{144.317}{448.063}{360.669}{ALIUSC1F8135}{\ALIUSFontLatoLight}{12.357}{also suggests something similarly. On the other hand, it is well-}
  \ALIUSPlacedTextContent{90.000}{465.103}{272.511}{ALIUSC1F8135}{\ALIUSFontLatoLight}{12.357}{known that psychedelic drugs can induce states of}
  \ALIUSPlacedTextContent{362.122}{465.103}{92.366}{ALIUSC1F8135}{\ALIUSFontLatoLightItalic}{12.357}{depersonalization}
  \ALIUSPlacedTextContent{454.165}{465.103}{54.218}{ALIUSC1F8135}{\ALIUSFontLatoLight}{12.357}{and more}
  \ALIUSPlacedTextContent{90.000}{482.383}{50.070}{ALIUSC1F8135}{\ALIUSFontLatoLight}{12.357}{radically,}
  \ALIUSPlacedTextContent{140.025}{482.383}{82.542}{ALIUSC1F8135}{\ALIUSFontLatoLightItalic}{12.357}{ego-dissolution,}
  \ALIUSPlacedTextContent{222.540}{482.383}{285.845}{ALIUSC1F8135}{\ALIUSFontLatoLight}{12.357}{in about 7\% of trip reports of high-dosage intakes of}
  \ALIUSPlacedTextContent{90.000}{499.663}{47.152}{ALIUSC1F8135}{\ALIUSFontLatoLightItalic}{12.357}{Psilocybe}
  \ALIUSPlacedTextContent{137.174}{499.663}{371.239}{ALIUSC1F8135}{\ALIUSFontLatoLight}{12.357}{mushrooms, LSD, Salvia, DMT, 5-MeO-DMT, ayahuasca and}
  \ALIUSPlacedTextContent{90.000}{516.943}{418.343}{ALIUSC1F8135}{\ALIUSFontLatoLight}{12.357}{ketamine (Millière, 2017). Despite being rare, these experiences have been}
  \ALIUSPlacedTextContent{90.000}{534.223}{316.563}{ALIUSC1F8135}{\ALIUSFontLatoLight}{12.357}{interpreted in relation to the Buddhist concept of no-self (}
  \ALIUSPlacedTextContent{406.568}{534.223}{45.637}{ALIUSC1F8135}{\ALIUSFontLatoLightItalic}{12.357}{anatman}
  \ALIUSPlacedTextContent{452.217}{534.223}{56.161}{ALIUSC1F8135}{\ALIUSFontLatoLight}{12.357}{) by some}
  \ALIUSPlacedTextContent{90.000}{551.503}{418.432}{ALIUSC1F8135}{\ALIUSFontLatoLight}{12.357}{early advocates of the use of psychedelics such as Timothy Leary or Aldous}
  \ALIUSPlacedTextContent{90.000}{568.783}{418.386}{ALIUSC1F8135}{\ALIUSFontLatoLight}{12.357}{Huxley (Huxley, 1999; Leary et al., 1964). What role do altered states of}
  \ALIUSPlacedTextContent{90.000}{586.063}{418.426}{ALIUSC1F8135}{\ALIUSFontLatoLight}{12.357}{consciousness play in Indian meditative practices? Do you think that we can}
  \ALIUSPlacedTextContent{90.000}{603.103}{418.480}{ALIUSC1F8135}{\ALIUSFontLatoLight}{12.357}{interpret the lack of self in conscious experiences under psychedelics in the}
  \ALIUSPlacedTextContent{90.000}{620.383}{163.272}{ALIUSC1F8135}{\ALIUSFontLatoLight}{12.357}{light of Buddhist philosophy?}
  \ALIUSPlacedTextContent{90.000}{649.538}{418.225}{ALIUSC000000}{\ALIUSFontCormorantRegular}{13.920}{I think it is better to describe the altered states of consciousness as}
  \ALIUSPlacedTextContent{90.000}{668.978}{418.295}{ALIUSC000000}{\ALIUSFontCormorantRegular}{13.920}{transformative experiences. These experiences, at least in the Buddhist}
  \ALIUSPlacedTextContent{90.000}{688.658}{418.310}{ALIUSC000000}{\ALIUSFontCormorantRegular}{13.920}{tradition are brought about by a variety of spiritual exercises. In keeping with}
  \ALIUSPlacedTextContent{90.000}{708.098}{292.464}{ALIUSC000000}{\ALIUSFontCormorantRegular}{13.920}{the general Buddhist tradition Vasubandhu in the}
  \ALIUSPlacedTextContent{385.414}{708.098}{119.526}{ALIUSC000000}{\ALIUSFontCormorantItalic}{13.920}{Abhidharmakośabhasya}
  \ALIUSPlacedTextContent{90.000}{727.538}{300.909}{ALIUSC000000}{\ALIUSFontCormorantRegular}{13.920}{notes that the spiritual path is an integrated system of}
  \ALIUSPlacedTextContent{391.270}{727.538}{17.053}{ALIUSC000000}{\ALIUSFontCormorantItalic}{13.920}{śīla}
  \ALIUSPlacedTextContent{408.349}{727.538}{96.625}{ALIUSC000000}{\ALIUSFontCormorantRegular}{13.920}{(moral conduct)-}
  \ALIUSPlacedTextContent{90.000}{746.978}{41.735}{ALIUSC000000}{\ALIUSFontCormorantItalic}{13.920}{samādhi}
  \ALIUSPlacedTextContent{131.747}{746.978}{76.903}{ALIUSC000000}{\ALIUSFontCormorantRegular}{13.920}{(meditation)-}
  \ALIUSPlacedTextContent{208.676}{746.978}{32.427}{ALIUSC000000}{\ALIUSFontCormorantItalic}{13.920}{prajñā}
  \ALIUSPlacedTextContent{241.141}{746.978}{270.392}{ALIUSC000000}{\ALIUSFontCormorantRegular}{13.920}{(wisdom or internalization of philosophical}
  \ALIUSPlacedTextContent{90.000}{792.192}{120.165}{ALIUSC767171}{\ALIUSFontCormorantLight}{12.000}{ALIUS Bulletin n°5 (2021)}
  \ALIUSPlacedTextContent{384.252}{792.192}{120.728}{ALIUSC767171}{\ALIUSFontCormorantLight}{12.000}{aliusresearch.org/bulletin}
\end{tikzpicture}
\null
\newpage
% Page 8
\thispagestyle{empty}
\noindent\begin{tikzpicture}[remember picture,overlay,x=1bp,y=1bp,shift={(current page.north west)}]
  \ALIUSPlacedTextContent{90.000}{35.378}{386.261}{ALIUSC767171}{\ALIUSFontCormorantRegular}{13.920}{Monima Chadha – Transformative experiences in Indian philosophy}
  \ALIUSPlacedTextContent{494.640}{35.472}{13.056}{ALIUSC000000}{\ALIUSFontCormorantRegular}{12.000}{36}
  \ALIUSPlacedTextContent{105.600}{83.477}{14.160}{ALIUSC595959}{\ALIUSFontCormorantLight}{48.000}{\ALIUSPullQuoteOpen}
  \ALIUSPlacedTextContent{133.440}{83.477}{326.728}{ALIUSC000000}{\ALIUSFontLatoLight}{15.120}{I think it is better to describe the altered states of}
  \ALIUSPlacedTextContent{472.800}{97.248}{14.160}{ALIUSC595959}{\ALIUSFontCormorantLight}{48.000}{\ALIUSPullQuoteClose}
  \ALIUSPlacedTextContent{133.440}{101.477}{292.978}{ALIUSC000000}{\ALIUSFontLatoLight}{15.120}{consciousness as transformative experiences.}
  \ALIUSPlacedTextContent{90.000}{149.378}{418.240}{ALIUSC000000}{\ALIUSFontCormorantRegular}{13.920}{insights of the Abhidharma tradition). In the commentary that follows,}
  \ALIUSPlacedTextContent{90.000}{168.818}{418.193}{ALIUSC000000}{\ALIUSFontCormorantRegular}{13.920}{Vasubandhu writes that whoever desires to see the truths should first of all}
  \ALIUSPlacedTextContent{90.000}{188.498}{89.380}{ALIUSC000000}{\ALIUSFontCormorantRegular}{13.920}{guard morality (}
  \ALIUSPlacedTextContent{179.425}{188.498}{17.053}{ALIUSC000000}{\ALIUSFontCormorantItalic}{13.920}{śīla}
  \ALIUSPlacedTextContent{196.504}{188.498}{311.719}{ALIUSC000000}{\ALIUSFontCormorantRegular}{13.920}{), then they read the teachings on which the insight into}
  \ALIUSPlacedTextContent{90.000}{207.938}{418.342}{ALIUSC000000}{\ALIUSFontCormorantRegular}{13.920}{the truths depends, and listen to their meaning, having listened they reflect}
  \ALIUSPlacedTextContent{90.000}{227.378}{418.245}{ALIUSC000000}{\ALIUSFontCormorantRegular}{13.920}{on the teachings, and having reflected, they devote themselves to the}
  \ALIUSPlacedTextContent{90.000}{246.818}{418.213}{ALIUSC000000}{\ALIUSFontCormorantRegular}{13.920}{cultivation of meditation. You can interpret the lack of self in conscious}
  \ALIUSPlacedTextContent{90.000}{266.498}{418.378}{ALIUSC000000}{\ALIUSFontCormorantRegular}{13.920}{experiences under psychedelics in the light of Buddhist philosophy, but I}
  \ALIUSPlacedTextContent{90.000}{285.938}{418.222}{ALIUSC000000}{\ALIUSFontCormorantRegular}{13.920}{think that would be unfair to the tradition. The experience of no self in}
  \ALIUSPlacedTextContent{90.000}{305.378}{418.098}{ALIUSC000000}{\ALIUSFontCormorantRegular}{13.920}{meditation is not an isolated phenomenon, it is embedded within the}
  \ALIUSPlacedTextContent{90.000}{324.818}{418.216}{ALIUSC000000}{\ALIUSFontCormorantRegular}{13.920}{spiritual exercises considered as a whole, including knowledge of Buddhist}
  \ALIUSPlacedTextContent{90.000}{344.258}{124.494}{ALIUSC000000}{\ALIUSFontCormorantRegular}{13.920}{philosophical insights.}
  \ALIUSPlacedTextContent{105.840}{395.237}{14.160}{ALIUSC595959}{\ALIUSFontCormorantLight}{48.000}{\ALIUSPullQuoteOpen}
  \ALIUSPlacedTextContent{132.960}{395.237}{327.313}{ALIUSC000000}{\ALIUSFontLatoLight}{15.120}{The experience of no self in meditation is not an}
  \ALIUSPlacedTextContent{132.960}{413.237}{327.336}{ALIUSC000000}{\ALIUSFontLatoLight}{15.120}{isolated phenomenon, it is embedded within the}
  \ALIUSPlacedTextContent{132.960}{431.237}{327.472}{ALIUSC000000}{\ALIUSFontLatoLight}{15.120}{spiritual exercises considered as a whole, including}
  \ALIUSPlacedTextContent{462.000}{445.968}{14.160}{ALIUSC595959}{\ALIUSFontCormorantLight}{48.000}{\ALIUSPullQuoteClose}
  \ALIUSPlacedTextContent{132.960}{449.237}{293.459}{ALIUSC000000}{\ALIUSFontLatoLight}{15.120}{knowledge of Buddhist philosophical insights.}
  \ALIUSPlacedTextContent{90.000}{500.623}{418.377}{ALIUSC1F8135}{\ALIUSFontLatoLight}{12.357}{In the Buddhist traditions, the notion of no-self has both metaphysical and}
  \ALIUSPlacedTextContent{90.000}{517.903}{418.244}{ALIUSC1F8135}{\ALIUSFontLatoLight}{12.357}{practical aspects, as holding the self as a substance is considered not only as}
  \ALIUSPlacedTextContent{90.000}{535.183}{418.250}{ALIUSC1F8135}{\ALIUSFontLatoLight}{12.357}{a false view about the reality but also as a source of suffering. How do}
  \ALIUSPlacedTextContent{90.000}{552.463}{418.437}{ALIUSC1F8135}{\ALIUSFontLatoLight}{12.357}{Buddhist traditions explain the therapeutic effects that can be attributed to a}
  \ALIUSPlacedTextContent{90.000}{569.743}{418.572}{ALIUSC1F8135}{\ALIUSFontLatoLight}{12.357}{change in our views about the self? Would you think that some aspects of}
  \ALIUSPlacedTextContent{90.000}{587.023}{418.355}{ALIUSC1F8135}{\ALIUSFontLatoLight}{12.357}{Buddhist philosophy could be useful to anchor therapeutic effects of}
  \ALIUSPlacedTextContent{90.000}{604.303}{418.394}{ALIUSC1F8135}{\ALIUSFontLatoLight}{12.357}{psychedelics by offering an understanding of the mental states being}
  \ALIUSPlacedTextContent{90.000}{621.583}{191.965}{ALIUSC1F8135}{\ALIUSFontLatoLight}{12.357}{experienced and their significance?}
  \ALIUSPlacedTextContent{90.000}{650.498}{418.101}{ALIUSC000000}{\ALIUSFontCormorantRegular}{13.920}{Belief in a continuing self is the basis of our special concern for our own}
  \ALIUSPlacedTextContent{90.000}{670.178}{418.297}{ALIUSC000000}{\ALIUSFontCormorantRegular}{13.920}{future self. We do put away money in superannuation rather than giving to}
  \ALIUSPlacedTextContent{90.000}{689.618}{418.240}{ALIUSC000000}{\ALIUSFontCormorantRegular}{13.920}{charity? Because we care more about our future self than contemporary}
  \ALIUSPlacedTextContent{90.000}{709.058}{418.277}{ALIUSC000000}{\ALIUSFontCormorantRegular}{13.920}{others, some of whom are suffering. The Buddhists do recognize that it is}
  \ALIUSPlacedTextContent{90.000}{728.498}{418.291}{ALIUSC000000}{\ALIUSFontCormorantRegular}{13.920}{built-in precondition of our form of life that we have self-concern and special}
  \ALIUSPlacedTextContent{90.000}{748.178}{418.226}{ALIUSC000000}{\ALIUSFontCormorantRegular}{13.920}{concern for our loved ones. That is why the Buddhists do not recommend}
  \ALIUSPlacedTextContent{90.000}{792.192}{120.165}{ALIUSC767171}{\ALIUSFontCormorantLight}{12.000}{ALIUS Bulletin n°5 (2021)}
  \ALIUSPlacedTextContent{384.252}{792.192}{120.728}{ALIUSC767171}{\ALIUSFontCormorantLight}{12.000}{aliusresearch.org/bulletin}
\end{tikzpicture}
\null
\newpage
% Page 9
\thispagestyle{empty}
\noindent\begin{tikzpicture}[remember picture,overlay,x=1bp,y=1bp,shift={(current page.north west)}]
  \ALIUSPlacedTextContent{90.000}{35.378}{386.261}{ALIUSC767171}{\ALIUSFontCormorantRegular}{13.920}{Monima Chadha – Transformative experiences in Indian philosophy}
  \ALIUSPlacedTextContent{495.120}{35.472}{12.624}{ALIUSC000000}{\ALIUSFontCormorantRegular}{12.000}{37}
  \ALIUSPlacedTextContent{90.000}{72.098}{418.278}{ALIUSC000000}{\ALIUSFontCormorantRegular}{13.920}{giving up on this self-concern or special concern for our loved ones. Rather}
  \ALIUSPlacedTextContent{90.000}{91.538}{418.320}{ALIUSC000000}{\ALIUSFontCormorantRegular}{13.920}{they recommend extending similar concern to others. And they do not think}
  \ALIUSPlacedTextContent{90.000}{110.978}{418.274}{ALIUSC000000}{\ALIUSFontCormorantRegular}{13.920}{that such an extension comes easy to us given our human nature: it has to be}
  \ALIUSPlacedTextContent{90.000}{130.658}{418.287}{ALIUSC000000}{\ALIUSFontCormorantRegular}{13.920}{inculcated by extensive meditation practices. As Parfit puts it, the discovery}
  \ALIUSPlacedTextContent{90.000}{150.098}{418.217}{ALIUSC000000}{\ALIUSFontCormorantRegular}{13.920}{of no self is liberating and consoling! Given what I've said in response to the}
  \ALIUSPlacedTextContent{90.000}{169.538}{418.122}{ALIUSC000000}{\ALIUSFontCormorantRegular}{13.920}{last question, I don't think we can better understand the therapeutic effects}
  \ALIUSPlacedTextContent{90.000}{188.978}{279.064}{ALIUSC000000}{\ALIUSFontCormorantRegular}{13.920}{of psychedelics by turning to Buddhist philosophy.}
  \ALIUSPlacedTextContent{90.000}{220.783}{24.039}{ALIUSC1F8135}{\ALIUSFontLatoLight}{12.357}{The}
  \ALIUSPlacedTextContent{119.183}{220.783}{62.400}{ALIUSC1F8135}{\ALIUSFontLatoLightItalic}{12.357}{Abhidharma}
  \ALIUSPlacedTextContent{181.594}{220.783}{326.790}{ALIUSC1F8135}{\ALIUSFontLatoLight}{12.357}{-Buddhist metaphysics of persons has some important}
  \ALIUSPlacedTextContent{90.000}{238.063}{418.372}{ALIUSC1F8135}{\ALIUSFontLatoLight}{12.357}{similarities to a 4D ontology of temporal parts (Sider, 2001), in particular, to a}
  \ALIUSPlacedTextContent{90.000}{255.103}{418.342}{ALIUSC1F8135}{\ALIUSFontLatoLight}{12.357}{'stage-theory' in which persons are momentary beings that strictly speaking}
  \ALIUSPlacedTextContent{90.000}{272.383}{418.404}{ALIUSC1F8135}{\ALIUSFontLatoLight}{12.357}{don't persist through time (impermanence). The stage theory has important}
  \ALIUSPlacedTextContent{90.000}{289.663}{418.380}{ALIUSC1F8135}{\ALIUSFontLatoLight}{12.357}{theoretical advantages in dealing with classical puzzles of the metaphysics of}
  \ALIUSPlacedTextContent{90.000}{306.943}{418.459}{ALIUSC1F8135}{\ALIUSFontLatoLight}{12.357}{material objects. However, when the stage theory is applied to persons (Olson,}
  \ALIUSPlacedTextContent{90.000}{324.223}{418.406}{ALIUSC1F8135}{\ALIUSFontLatoLight}{12.357}{2007), it faces an immediate objection: diachronic features seem rather crucial}
  \ALIUSPlacedTextContent{90.000}{341.503}{418.407}{ALIUSC1F8135}{\ALIUSFontLatoLight}{12.357}{in personal identity. Even if not for metaphysical reasons, some practical}
  \ALIUSPlacedTextContent{90.000}{358.783}{186.561}{ALIUSC1F8135}{\ALIUSFontLatoLight}{12.357}{affairs seem essentially diachronic (}
  \ALIUSPlacedTextContent{276.660}{358.783}{20.053}{ALIUSC1F8135}{\ALIUSFontLatoLightItalic}{12.357}{e.g.,}
  \ALIUSPlacedTextContent{296.710}{358.783}{211.638}{ALIUSC1F8135}{\ALIUSFontLatoLight}{12.357}{norms of rationality, responsibility, and}
  \ALIUSPlacedTextContent{90.000}{376.063}{182.812}{ALIUSC1F8135}{\ALIUSFontLatoLight}{12.357}{regret, etc). How do you think the}
  \ALIUSPlacedTextContent{272.243}{376.063}{62.400}{ALIUSC1F8135}{\ALIUSFontLatoLightItalic}{12.357}{Abhidharma}
  \ALIUSPlacedTextContent{334.654}{376.063}{173.732}{ALIUSC1F8135}{\ALIUSFontLatoLight}{12.357}{theory of persons can deal with}
  \ALIUSPlacedTextContent{90.000}{393.103}{76.037}{ALIUSC1F8135}{\ALIUSFontLatoLight}{12.357}{this problem?}
  \ALIUSPlacedTextContent{90.000}{422.258}{418.218}{ALIUSC000000}{\ALIUSFontCormorantRegular}{13.920}{The central normative goal of Buddhism is to ameliorate suffering and that}
  \ALIUSPlacedTextContent{90.000}{441.938}{418.249}{ALIUSC000000}{\ALIUSFontCormorantRegular}{13.920}{guides its revision of descriptive metaphysics. The no-self and no-person}
  \ALIUSPlacedTextContent{90.000}{461.378}{418.093}{ALIUSC000000}{\ALIUSFontCormorantRegular}{13.920}{metaphysics aims to produce a better structure that is motivated by the}
  \ALIUSPlacedTextContent{90.000}{480.818}{418.282}{ALIUSC000000}{\ALIUSFontCormorantRegular}{13.920}{normative goal of eliminating, or at least reducing, suffering. Buddhist}
  \ALIUSPlacedTextContent{90.000}{500.258}{418.286}{ALIUSC000000}{\ALIUSFontCormorantRegular}{13.920}{revisionary metaphysics is not aimed at capturing the structure the world}
  \ALIUSPlacedTextContent{90.000}{519.698}{418.229}{ALIUSC000000}{\ALIUSFontCormorantRegular}{13.920}{really has, and a justification of our ordinary person-related practices. Rather}
  \ALIUSPlacedTextContent{90.000}{539.378}{418.252}{ALIUSC000000}{\ALIUSFontCormorantRegular}{13.920}{it aims at providing a structure that aims to reduce suffering. The revised}
  \ALIUSPlacedTextContent{90.000}{558.818}{418.343}{ALIUSC000000}{\ALIUSFontCormorantRegular}{13.920}{structure, in turn, entails a major reconsideration of our ordinary everyday}
  \ALIUSPlacedTextContent{90.000}{578.258}{418.192}{ALIUSC000000}{\ALIUSFontCormorantRegular}{13.920}{person-related concerns and practices and interpersonal attitudes, such as}
  \ALIUSPlacedTextContent{90.000}{597.698}{415.184}{ALIUSC000000}{\ALIUSFontCormorantRegular}{13.920}{moral responsibility, praise and blame, compensation, and social treatment.}
  \ALIUSPlacedTextContent{90.000}{629.503}{418.280}{ALIUSC1F8135}{\ALIUSFontLatoLight}{12.357}{ALIUS is most interested in the diversity of consciousness. This involves both}
  \ALIUSPlacedTextContent{90.000}{646.783}{418.380}{ALIUSC1F8135}{\ALIUSFontLatoLight}{12.357}{discussing the diversity of conscious states and the diversity of disciplines or}
  \ALIUSPlacedTextContent{90.000}{663.823}{418.394}{ALIUSC1F8135}{\ALIUSFontLatoLight}{12.357}{cultural outlooks that can inform a scientific understanding of the nature of}
  \ALIUSPlacedTextContent{90.000}{681.103}{418.463}{ALIUSC1F8135}{\ALIUSFontLatoLight}{12.357}{consciousness. What are the challenges according to you in developing a}
  \ALIUSPlacedTextContent{90.000}{698.383}{418.521}{ALIUSC1F8135}{\ALIUSFontLatoLight}{12.357}{cross-cultural discussion about consciousness and the mind? What is your}
  \ALIUSPlacedTextContent{90.000}{715.663}{226.397}{ALIUSC1F8135}{\ALIUSFontLatoLight}{12.357}{approach to overcome these challenges?}
  \ALIUSPlacedTextContent{90.000}{750.098}{418.260}{ALIUSC000000}{\ALIUSFontCormorantRegular}{13.920}{One of the most important challenges is to draw the attention of mainstream}
  \ALIUSPlacedTextContent{90.000}{792.192}{120.165}{ALIUSC767171}{\ALIUSFontCormorantLight}{12.000}{ALIUS Bulletin n°5 (2021)}
  \ALIUSPlacedTextContent{384.252}{792.192}{120.728}{ALIUSC767171}{\ALIUSFontCormorantLight}{12.000}{aliusresearch.org/bulletin}
\end{tikzpicture}
\null
\newpage
% Page 10
\thispagestyle{empty}
\noindent\begin{tikzpicture}[remember picture,overlay,x=1bp,y=1bp,shift={(current page.north west)}]
  \fill[ALIUSC0000FF] (204.000bp,-313.440bp) rectangle ++(197.280bp,-0.720bp);
  \fill[ALIUSC0000FF] (149.520bp,-372.720bp) rectangle ++(217.440bp,-0.720bp);
  \fill[ALIUSC0000FF] (90.000bp,-561.600bp) rectangle ++(187.200bp,-0.720bp);
  \fill[ALIUSC0000FF] (229.440bp,-671.760bp) rectangle ++(206.640bp,-0.720bp);
  \ALIUSPlacedTextContent{90.000}{35.378}{386.261}{ALIUSC767171}{\ALIUSFontCormorantRegular}{13.920}{Monima Chadha – Transformative experiences in Indian philosophy}
  \ALIUSPlacedTextContent{494.400}{35.472}{13.344}{ALIUSC000000}{\ALIUSFontCormorantRegular}{12.000}{38}
  \ALIUSPlacedTextContent{90.000}{72.098}{418.213}{ALIUSC000000}{\ALIUSFontCormorantRegular}{13.920}{philosophers of mind and consciousness to pay serious attention to the}
  \ALIUSPlacedTextContent{90.000}{91.538}{418.276}{ALIUSC000000}{\ALIUSFontCormorantRegular}{13.920}{classical Indian tradition and what it has to offer. There is not just the}
  \ALIUSPlacedTextContent{90.000}{110.978}{418.199}{ALIUSC000000}{\ALIUSFontCormorantRegular}{13.920}{language barrier but also the misconception that Indian philosophy is mostly}
  \ALIUSPlacedTextContent{90.000}{130.658}{418.285}{ALIUSC000000}{\ALIUSFontCormorantRegular}{13.920}{mystical mumbo-jumbo. Correcting that misconception is a hard task. So my}
  \ALIUSPlacedTextContent{90.000}{150.098}{418.310}{ALIUSC000000}{\ALIUSFontCormorantRegular}{13.920}{approach has been trying to publish in mainstream journals rather than}
  \ALIUSPlacedTextContent{90.000}{169.538}{418.424}{ALIUSC000000}{\ALIUSFontCormorantRegular}{13.920}{specialist Indian philosophy journals so that these materials have a chance of}
  \ALIUSPlacedTextContent{90.000}{188.978}{305.244}{ALIUSC000000}{\ALIUSFontCormorantRegular}{13.920}{broader uptake. It is hard, but I think it is worth doing.}
  \ALIUSPlacedTextContent{263.666}{238.021}{71.454}{ALIUSC000000}{\ALIUSFontLatoLight}{13.920}{References}
  \ALIUSPlacedTextContent{126.000}{269.065}{382.033}{ALIUSC000000}{\ALIUSFontCormorantRegular}{12.960}{Chadha, M. (2019). Abhidharma Panprotopsychist Metaphysics of}
  \ALIUSPlacedTextContent{90.000}{284.665}{178.983}{ALIUSC000000}{\ALIUSFontCormorantRegular}{12.960}{Consciousness. In W. Seager (Ed.),}
  \ALIUSPlacedTextContent{269.696}{284.665}{193.392}{ALIUSC000000}{\ALIUSFontCormorantItalic}{12.960}{The Routledge Handbook of Panpsychism}
  \ALIUSPlacedTextContent{463.125}{284.665}{44.908}{ALIUSC000000}{\ALIUSFontCormorantRegular}{12.960}{(1st ed.,}
  \ALIUSPlacedTextContent{90.000}{300.505}{113.873}{ALIUSC000000}{\ALIUSFontCormorantRegular}{12.960}{pp. 25–35). Routledge.}
  \ALIUSPlacedTextContent{203.903}{300.505}{197.254}{ALIUSC0000FF}{\ALIUSFontCormorantRegular}{12.960}{https://doi.org/10.4324/9781315717708-3}
  \ALIUSPlacedTextContent{126.000}{328.105}{381.984}{ALIUSC000000}{\ALIUSFontCormorantRegular}{12.960}{Chalmers, D. (2007). The Hard Problem of Consciousness. In M. Velmans \&}
  \ALIUSPlacedTextContent{90.000}{343.945}{98.729}{ALIUSC000000}{\ALIUSFontCormorantRegular}{12.960}{S. Schneider (Eds.),}
  \ALIUSPlacedTextContent{188.093}{343.945}{196.971}{ALIUSC000000}{\ALIUSFontCormorantItalic}{12.960}{The Blackwell Companion to Consciousness}
  \ALIUSPlacedTextContent{385.082}{343.945}{122.967}{ALIUSC000000}{\ALIUSFontCormorantRegular}{12.960}{(pp. 223–235). Blackwell}
  \ALIUSPlacedTextContent{90.000}{359.785}{59.534}{ALIUSC000000}{\ALIUSFontCormorantRegular}{12.960}{Publishing.}
  \ALIUSPlacedTextContent{149.538}{359.785}{217.333}{ALIUSC0000FF}{\ALIUSFontCormorantRegular}{12.960}{https://doi.org/10.1002/9780470751466.ch18}
  \ALIUSPlacedTextContent{126.000}{391.225}{327.476}{ALIUSC000000}{\ALIUSFontCormorantRegular}{12.960}{Frankish, K. (2016). Illusionism as a theory of consciousness.}
  \ALIUSPlacedTextContent{456.441}{391.225}{51.568}{ALIUSC000000}{\ALIUSFontCormorantItalic}{12.960}{Journal of}
  \ALIUSPlacedTextContent{90.000}{406.825}{100.246}{ALIUSC000000}{\ALIUSFontCormorantItalic}{12.960}{Consciousness Studies}
  \ALIUSPlacedTextContent{190.254}{406.825}{5.919}{ALIUSC000000}{\ALIUSFontCormorantRegular}{12.960}{,}
  \ALIUSPlacedTextContent{196.182}{406.825}{9.151}{ALIUSC000000}{\ALIUSFontCormorantItalic}{12.960}{23}
  \ALIUSPlacedTextContent{205.346}{406.825}{70.852}{ALIUSC000000}{\ALIUSFontCormorantRegular}{12.960}{(11–12), 11–39.}
  \ALIUSPlacedTextContent{126.000}{438.505}{382.035}{ALIUSC000000}{\ALIUSFontCormorantRegular}{12.960}{Goff, P. (2017). Consciousness and Fundamental Reality. Oxford: Oxford}
  \ALIUSPlacedTextContent{90.000}{454.105}{87.530}{ALIUSC000000}{\ALIUSFontCormorantRegular}{12.960}{University Press.}
  \ALIUSPlacedTextContent{126.000}{485.545}{94.185}{ALIUSC000000}{\ALIUSFontCormorantRegular}{12.960}{Huxley, A. (1999).}
  \ALIUSPlacedTextContent{220.788}{485.545}{287.297}{ALIUSC000000}{\ALIUSFontCormorantItalic}{12.960}{Moksha: Aldous Huxley's classic writings on psychedelics and}
  \ALIUSPlacedTextContent{90.000}{501.385}{109.858}{ALIUSC000000}{\ALIUSFontCormorantItalic}{12.960}{the visionary experience}
  \ALIUSPlacedTextContent{199.926}{501.385}{112.347}{ALIUSC000000}{\ALIUSFontCormorantRegular}{12.960}{. Simon and Schuster.}
  \ALIUSPlacedTextContent{126.000}{532.825}{248.534}{ALIUSC000000}{\ALIUSFontCormorantRegular}{12.960}{Kriegel, U. (2019). The Value of Consciousness.}
  \ALIUSPlacedTextContent{376.066}{532.825}{39.380}{ALIUSC000000}{\ALIUSFontCormorantItalic}{12.960}{Analysis}
  \ALIUSPlacedTextContent{415.454}{532.825}{92.574}{ALIUSC000000}{\ALIUSFontCormorantRegular}{12.960}{, 79, (3), 503–520.}
  \ALIUSPlacedTextContent{90.000}{548.665}{187.289}{ALIUSC0000FF}{\ALIUSFontCormorantRegular}{12.960}{https://doi.org/10.1093/analys/anz045}
  \ALIUSPlacedTextContent{126.000}{580.105}{206.974}{ALIUSC000000}{\ALIUSFontCormorantRegular}{12.960}{Leary, T., Metzner, R., \& Dass, R. (1964).}
  \ALIUSPlacedTextContent{332.877}{580.105}{175.105}{ALIUSC000000}{\ALIUSFontCormorantItalic}{12.960}{The psychedelic experience: A manual}
  \ALIUSPlacedTextContent{90.000}{595.945}{243.110}{ALIUSC000000}{\ALIUSFontCormorantItalic}{12.960}{based on the Tibetan book of the dead (Vol. 1)}
  \ALIUSPlacedTextContent{333.103}{595.945}{174.930}{ALIUSC000000}{\ALIUSFontCormorantRegular}{12.960}{. Carol Publishing Corporation.}
  \ALIUSPlacedTextContent{126.000}{627.385}{381.994}{ALIUSC000000}{\ALIUSFontCormorantRegular}{12.960}{Millière, R. (2017). Looking for the Self: Phenomenology, Neurophysiology}
  \ALIUSPlacedTextContent{90.000}{642.985}{353.387}{ALIUSC000000}{\ALIUSFontCormorantRegular}{12.960}{and Philosophical Significance of Drug-induced Ego Dissolution.}
  \ALIUSPlacedTextContent{447.140}{642.985}{60.873}{ALIUSC000000}{\ALIUSFontCormorantItalic}{12.960}{Frontiers in}
  \ALIUSPlacedTextContent{90.000}{658.825}{97.896}{ALIUSC000000}{\ALIUSFontCormorantItalic}{12.960}{Human Neuroscience}
  \ALIUSPlacedTextContent{187.914}{658.825}{5.919}{ALIUSC000000}{\ALIUSFontCormorantRegular}{12.960}{,}
  \ALIUSPlacedTextContent{193.842}{658.825}{7.554}{ALIUSC000000}{\ALIUSFontCormorantItalic}{12.960}{11}
  \ALIUSPlacedTextContent{201.408}{658.825}{27.927}{ALIUSC000000}{\ALIUSFontCormorantRegular}{12.960}{, 245.}
  \ALIUSPlacedTextContent{229.344}{658.825}{206.615}{ALIUSC0000FF}{\ALIUSFontCormorantRegular}{12.960}{https://doi.org/10.3389/fnhum.2017.00245}
  \ALIUSPlacedTextContent{126.000}{690.265}{85.693}{ALIUSC000000}{\ALIUSFontCormorantRegular}{12.960}{Olson, E. (2007).}
  \ALIUSPlacedTextContent{211.546}{690.265}{207.995}{ALIUSC000000}{\ALIUSFontCormorantItalic}{12.960}{What Are We? A Study in Personal Ontology}
  \ALIUSPlacedTextContent{419.613}{690.265}{88.422}{ALIUSC000000}{\ALIUSFontCormorantRegular}{12.960}{. Oxford: Oxford}
  \ALIUSPlacedTextContent{90.000}{706.105}{87.530}{ALIUSC000000}{\ALIUSFontCormorantRegular}{12.960}{University Press.}
  \ALIUSPlacedTextContent{198.000}{706.105}{8.635}{ALIUSC000000}{\ALIUSFontCormorantRegular}{12.960}{.}
  \ALIUSPlacedTextContent{126.000}{737.545}{86.074}{ALIUSC000000}{\ALIUSFontCormorantRegular}{12.960}{Sider, T. (2001).}
  \ALIUSPlacedTextContent{214.162}{737.545}{293.802}{ALIUSC000000}{\ALIUSFontCormorantItalic}{12.960}{Four-Dimensionalism: An Ontology of Persistence and Time.}
  \ALIUSPlacedTextContent{90.000}{753.385}{170.663}{ALIUSC000000}{\ALIUSFontCormorantRegular}{12.960}{Oxford: Oxford University Press.}
  \ALIUSPlacedTextContent{90.000}{792.192}{120.165}{ALIUSC767171}{\ALIUSFontCormorantLight}{12.000}{ALIUS Bulletin n°5 (2021)}
  \ALIUSPlacedTextContent{384.252}{792.192}{120.728}{ALIUSC767171}{\ALIUSFontCormorantLight}{12.000}{aliusresearch.org/bulletin}
\end{tikzpicture}
\null
\ifdefined\ALIUSIssueBuild
\clearpage
\expandafter\endinput
\fi
\end{document}
